% !TEX program = xelatex
%=========================================================================
%  基于 AIGC 的课程知识图谱智能构建与学习导航系统 · 系统使用说明文档
%
%  由 md2tex.py 自动生成，请勿直接编辑本文件——
%  修订请改 docs/系统使用说明文档.md 后重新运行：
%      python md2tex.py docs/系统使用说明文档.md docs/系统使用说明文档.tex
%
%  ★ 编译方式：必须使用 XeLaTeX（Overleaf：Menu → Compiler → XeLaTeX）
%    默认的 pdfLaTeX 无法处理中文，会直接编译失败。
%=========================================================================
\documentclass[UTF8,a4paper,11pt]{ctexrep}

\usepackage{geometry}
\geometry{left=2.4cm,right=2.4cm,top=2.6cm,bottom=2.6cm}

\usepackage{xcolor}
\definecolor{notebar}{HTML}{4F6EF7}    % 提示框左侧竖条（蓝）
\definecolor{notebg}{HTML}{F2F5FF}     % 提示框底色
\definecolor{todobar}{HTML}{E6A23C}    % 待办框左侧竖条（橙）
\definecolor{todobg}{HTML}{FFF8EC}     % 待办框底色
\definecolor{codeframe}{HTML}{D0D7E2}  % 代码块边框

\usepackage{amsmath}
\usepackage{amssymb}
\usepackage{booktabs}
\usepackage{tabularx}
\usepackage{array}
\usepackage{enumitem}
\usepackage{fvextra}       % 提供可换行的 Verbatim
\usepackage{newunicodechar}
\usepackage{needspace}

% ★ 等宽字体：必须换成一个含制表符（─│┌┐└┘）的字体，默认的 Latin Modern Mono 没有这些字形。
%   Scale=0.83 是刻意的：让西文宽度恰为中文字符的一半，
%   这样文档里的架构图/目录树才能对齐（原文按「中文占 2 个西文字符宽」书写）。
%   若报「字体找不到」，把下面一行换成：\setmonofont{FreeMono}[Scale=0.83]
\setmonofont{DejaVu Sans Mono}[Scale=0.83]

% 箭头、比较号等：交由数学模式渲染，避免西文字体缺字
\newunicodechar{→}{\ensuremath{\rightarrow}}
\newunicodechar{↓}{\ensuremath{\downarrow}}
\newunicodechar{≥}{\ensuremath{\geq}}
\newunicodechar{≤}{\ensuremath{\leq}}
\newunicodechar{⊆}{\ensuremath{\subseteq}}
\newunicodechar{▼}{\ensuremath{\blacktriangledown}}
\newunicodechar{×}{\ensuremath{\times}}
\newunicodechar{·}{\textperiodcentered}
\newunicodechar{…}{\textellipsis}

\usepackage{hyperref}
\hypersetup{hidelinks,bookmarksnumbered,bookmarksopen}
% tcolorbox 官方建议在 hyperref 之后载入
\usepackage{tcolorbox}

% ---- 提示框 ----
% 注意：两个框都刻意做成「不可分页」。
% 可分页盒子（breakable）与 tabularx 组合会出现排版异常，
% 而本文档的提示框均短于一页，不需要分页。
\newtcolorbox{notebox}{
  colback=notebg, colframe=notebar, boxrule=0pt, leftrule=3pt,
  arc=1pt, left=8pt, right=6pt, top=6pt, bottom=6pt, fontupper=\small,
  before skip=8pt, after skip=8pt
}
\newtcolorbox{todobox}{
  colback=todobg, colframe=todobar, boxrule=0pt, leftrule=3pt,
  arc=1pt, left=8pt, right=6pt, top=6pt, bottom=6pt, fontupper=\small,
  before skip=8pt, after skip=8pt
}

% ---- 占位符标记（对应 Markdown 中的 emoji 约定）----
\newcommand{\shotmark}{\textcolor{todobar}{\textbf{[截图待补]}}}
\newcommand{\todomark}{\textcolor{todobar}{\textbf{[数据待补]}}}
\newcommand{\verifymark}{\textcolor{todobar}{\textbf{[待确认]}}}
\newcommand{\warnmark}{\textcolor{todobar}{\textbf{[注意]}}}

% ---- 版面微调 ----
\setlength{\parskip}{0.35em}
\linespread{1.12}
\renewcommand{\arraystretch}{1.28}
\setlist[itemize]{leftmargin=1.2em,itemsep=2pt,topsep=3pt,parsep=0pt}
\setlist[enumerate]{leftmargin=1.6em,itemsep=2pt,topsep=3pt,parsep=0pt}

\begin{document}

\begin{center}
{\Huge\bfseries 基于 AIGC 的课程知识图谱智能构建与学习导航系统\par}
\end{center}


\begin{center}
{\LARGE 系统使用说明文档\par}
\end{center}
\vspace{1em}

\begin{center}
\small
\begin{tabularx}{\linewidth}{@{}XX@{}}
\toprule
项 & 内容 \\
\midrule
文档版本 & v2.0 \\
对应系统版本 & 见 \texttt{backend/app/core/config.py} 的 \texttt{VERSION} \\
发布日期 & 待填 \\
适用对象 & 部署人员、教师、学生、管理员 \\
配套材料 & 《知识抽取准确率测试报告》《问答测试集》《提示词工程完整记录》《项目详细方案》 \\
\bottomrule
\end{tabularx}
\end{center}

\subsection*{修订记录}

\begin{center}
\small
\begin{tabularx}{\linewidth}{@{}XXXX@{}}
\toprule
版本 & 日期 & 修订内容 & 修订人 \\
\midrule
v1.0 & — & 初版：安装部署、操作流程、常见问题 & — \\
v2.0 & — & 按竞赛交付要求重构为七章结构；新增快速启动、部署验收、赛题指标自检、系统测试、术语表；按角色拆分 FAQ；修正 RAG 与路径算法表述 & — \\
\bottomrule
\end{tabularx}
\end{center}
\begin{todobox}

\textbf{\warnmark{} 提交前必读：文档交付状态说明}\par

本文档中凡标记 \textbf{\todomark{}} 的内容为\textbf{待补占位项}，需团队在提交前用\textbf{真实截图}与\textbf{实测数据}替换。\textbf{请勿直接提交含占位符的版本。}

\begin{center}
\small
\begin{tabularx}{\linewidth}{@{}XXXX@{}}
\toprule
占位类型 & 标记 & 数量 & 说明 \\
\midrule
界面截图 & \shotmark{} & 6 处 & 需在真实运行的系统上截取，不可使用示意图代替 \\
实测数值 & \todomark{} & 3 张表 & 性能指标、功能测试结果、多课程验证记录 \\
待确认信息 & \verifymark{} & 若干 & 需团队核对后填写（如发布日期、系统版本号、仓库地址） \\
\bottomrule
\end{tabularx}
\end{center}
\textbf{凡是数值一律标注来源。} 来自已有测试报告的数字标注了出处文件名；尚无实测记录的表格留空并注明采集方法，\textbf{不用预期值或估计值填充}。

\textbf{关于截图的取舍}：本文档只保留 \textbf{6 张}截图，选取标准是「\textbf{是否直接支撑赛题的评分项}」——登录页、导航栏、成员管理、题库、学情监测等界面虽有截图会更好看，但它们不对应任何硬指标，故只作文字说明。保留的 6 张与赛题要求的对应关系见 附录 E。

\begin{center}
\small
\begin{tabularx}{\linewidth}{@{}XXX@{}}
\toprule
编号 & 截图 & 支撑的赛题要求 \\
\midrule
3-1 & 文档上传与处理结果 & 支持 ≥2 种文档格式；知识点 ≥20 个 \\
3-2 & 知识图谱浏览 & 图谱可视化与交互；关系类型 ≥3 种；知识点详情展示 \\
3-3 & 图谱编辑（前后对比） & 教师端「手动修正图谱」 \\
3-4 & 文档阅读与知识点联动 & 知识点可追溯至原文（提交材料第 5 项） \\
3-5 & 学习路径推荐 & 个性化学习路径推荐；路径可视化 \\
3-6 & 智能问答与引用来源 & 基于课程知识图谱的智能问答（RAG） \\
\bottomrule
\end{tabularx}
\end{center}
\end{todobox}


\subsection*{按角色快速阅读}

\begin{center}
\small
\begin{tabularx}{\linewidth}{@{}XX@{}}
\toprule
如果你是… & 建议阅读 \\
\midrule
部署人员 & 第二章（安装部署）→ 第五章（指标验证）→ 第七章（运维） \\
教师 & 3.1（导航）→ 3.2（教师端）→ 4.1–4.2（抽取与图谱原理）→ 6.2（教师 FAQ） \\
学生 & 3.1（导航）→ 3.3（学生端）→ 4.3–4.4（路径与问答原理）→ 6.3（学生 FAQ） \\
管理员 & 3.4（管理员端）→ 7.2–7.3（账号与日志）→ 6.4（管理员 FAQ） \\
评委 & 1.3（功能特色）→ 第三章（功能与截图）→ 第五章（指标验证）→ 附录 E（赛题要求映射表） \\
\bottomrule
\end{tabularx}
\end{center}
\tableofcontents
\clearpage


\chapter{系统概述}


\section{系统简介与建设目标}

在高校教学过程中，课程内容往往以教材、课件、教案等\textbf{非结构化文档}形式存在，知识点分散、缺乏清晰的关联关系。传统课程知识图谱的构建高度依赖教师或专家手工整理，工作量大、周期长，难以推广。

本系统利用 AIGC 技术自动完成 \textbf{文档解析 → 知识点抽取 → 关系构建 → 图谱生成}，把知识图谱的构建门槛从「专家标注」降低到「上传一份文档」，并在此基础上提供个性化学习路径推荐与基于 RAG 的智能问答。

\textbf{建设目标}：

\begin{enumerate}
\item 让教师无需具备图谱建模知识，也能为学生提供结构化的课程知识资源；
\item 让学生从「线性阅读教材」转向「按知识关联学习」；
\item 让学习路径从「教师统一安排」转向「依据个人掌握情况动态推荐」。
\end{enumerate}

\section{适用对象与使用场景}

\begin{center}
\small
\begin{tabularx}{\linewidth}{@{}XX@{}}
\toprule
角色 & 典型使用场景 \\
\midrule
高校教师 & 上传本课程教材/课件 → 自动生成知识图谱 → 人工校对修正 → 分享给学生 → 通过学情监测了解学习进度 \\
高校学生 & 加入课程 → 浏览知识图谱理解课程结构 → 标记已掌握知识点 → 按推荐路径学习 → 随时提问 \\
教学管理者 & 多课程统一管理、平台资源治理、运行状态监控 \\
\bottomrule
\end{tabularx}
\end{center}
\textbf{不适用场景（边界说明）}：

\begin{itemize}
\item 系统面向\textbf{单门课程内的知识组织}，明确\textbf{不提供跨课程知识融合}（不同课程的同名知识点相互独立）；
\item 抽取质量受原文结构影响，\textbf{扫描版 PDF（无文本层）无法处理}；
\item 系统依赖外部大模型 API，\textbf{必须联网}，不适用于完全离线环境。
\end{itemize}

\section{系统核心功能与特色}

以下五项是本系统相对手工建库方式的实质差异，也是竞赛演示中应当重点展示的部分。


\subsection{特色一：从文档到图谱的自动化构建}

教师上传 PDF / DOCX / TXT / MD 文档后，系统自动完成解析、分块、实体抽取、关系识别与入图，无需任何标注工作。

\begin{itemize}
\item \textbf{对教师的价值}：原本需要数天手工整理的知识结构，压缩到一次上传；
\item \textbf{可验证的证据}：见 3.2.3 上传文档 与 5.3 准确率测试。
\end{itemize}

\subsection{特色二：知识点在原文中的可追溯定位}

知识图谱中的每个知识点都可以回到它被抽取出来的\textbf{原始文档片段}。

\begin{itemize}
\item \textbf{对教师的价值}：可以核对抽取是否准确，判断是模型错误还是原文本身表述模糊；
\item \textbf{对学生的价值}：阅读文档时可看到当前段落对应的知识点，实现「读文献 + 看图谱」联动；
\item \textbf{对应功能}：见 3.2.9 文档在线阅读。
\end{itemize}

\subsection{特色三：基于掌握情况与前置关系的路径推荐}

系统依据图谱中的前置关系和学生的实际掌握情况，动态计算下一步应学习的知识点。

\begin{itemize}
\item \textbf{对学生的价值}：不必自己判断「接下来该学什么」，避免在不具备前置知识时贸然学习后续内容；
\item \textbf{算法约定与边界条件}：见 4.3 个性化学习路径推荐。
\end{itemize}

\subsection{特色四：基于课程内容的 RAG 问答}

学生提问后，系统先检索课程知识图谱中的相关知识点，再交由大模型生成回答，并\textbf{给出本次回答所依据的知识点引用来源}。

\begin{itemize}
\item \textbf{对学生的价值}：回答围绕本课程内容展开，并可通过引用来源自行核对；
\item \textbf{技术边界说明}：见 4.4 RAG 智能问答（含准确性的客观描述）。
\end{itemize}

\subsection{特色五：教师可校对、可干预}

自动生成的结果并非最终结果。教师可增删节点、修改知识点描述、建立关系，也可对知识点图谱结构进行人工修正。

\begin{itemize}
\item \textbf{对应功能}：见 3.2.4 查看与编辑知识图谱。
\end{itemize}

\section{系统架构与技术环境}


\subsection{总体架构}

\begin{Verbatim}[fontsize=\small,frame=single,framesep=4pt,rulecolor=\color{codeframe}]
┌─────────────────────────────────────────────────────┐
│  前端  Vue 3 + Vite + Pinia + Element Plus          │
│        知识图谱渲染：AntV G6 v5                       │
└────────────────────┬────────────────────────────────┘
                     │ HTTP（开发环境经 Vite 代理）
┌────────────────────▼────────────────────────────────┐
│  后端  FastAPI + Uvicorn                             │
│        分层：api（路由）→ services（业务）→ core（基础）│
└──────┬──────────────────────────────┬───────────────┘
       │                              │
┌──────▼──────────┐          ┌────────▼──────────────┐
│  Neo4j 5.x      │          │  SQLite (app.db)      │
│  知识图谱本体    │          │  用户/课程/文档/题库    │
│  节点与关系      │          │  学习记录/向量          │
└─────────────────┘          └───────────────────────┘
       │
┌──────▼──────────────────────────────────────────────┐
│  外部服务：DeepSeek（抽取与生成）                      │
│           SiliconFlow bge-m3（向量化）                │
└─────────────────────────────────────────────────────┘
\end{Verbatim}


\subsection{技术选型对照（赛题推荐路线）}

\begin{center}
\small
\begin{tabularx}{\linewidth}{@{}XXXX@{}}
\toprule
类别 & 赛题建议 & 本系统采用 & 说明 \\
\midrule
大模型 API & DeepSeek、阿里通义千问 & \textbf{DeepSeek} \texttt{deepseek-chat} & OpenAI 兼容接口 \\
图数据库 & Neo4j 社区版 / AuraDB & \textbf{Neo4j 5.x 社区版} & 本地部署 \\
可视化库 & ECharts Graph、AntV G6、D3.js、Vis.js & \textbf{AntV G6 v5} & — \\
后端框架 & Flask / FastAPI / Django & \textbf{FastAPI} & — \\
前端框架 & Vue / React / Streamlit & \textbf{Vue 3 + Vite} & — \\
文档解析 & PyPDF2、pdfplumber、python-docx & \textbf{三者均使用} & 按格式分别调用 \\
RAG 框架 & LangChain、LlamaIndex（可选） & \textbf{自研检索链路} & 见 4.4 说明 \\
\bottomrule
\end{tabularx}
\end{center}
\begin{notebox}
\textbf{关于未使用 RAG 框架}：赛题注明 RAG 框架为可选。本项目采用自研检索链路（向量化 → 余弦相似度排序 → 拼接上下文），原因是在当前文档规模下（单文档数十至数百个知识点）无需引入 ANN 索引，且自研链路便于在《论文算法复现实验》中与论文算法做对照实验。

\end{notebox}


\subsection{技术环境}

\begin{center}
\small
\begin{tabularx}{\linewidth}{@{}XX@{}}
\toprule
项 & 版本/要求 \\
\midrule
Python & ≥ 3.10 \\
Node.js & ≥ 18 \\
Neo4j & 5.x \\
操作系统 & Windows / macOS / Linux \\
浏览器 & Chrome / Edge 最新版（图谱渲染依赖 Canvas 与 WebGL） \\
GPU & 不需要 \\
\bottomrule
\end{tabularx}
\end{center}

\section{用户角色与权限说明}

系统分三种角色，权限按\textbf{三档}授予。

\begin{center}
\small
\begin{tabularx}{\linewidth}{@{}XXX@{}}
\toprule
角色 & 登录后落点 & 可访问范围 \\
\midrule
教师 & 课程中心 → 我的课程 & 本人创建或被邀请协作的课程 \\
学生 & 课程中心 → 我的课程 & 已加入并审核通过的课程 \\
管理员 & 管理工作台 & 全平台（治理视角，不参与教学） \\
\bottomrule
\end{tabularx}
\end{center}

\subsection{三档权限}

\begin{center}
\small
\begin{tabularx}{\linewidth}{@{}XXXX@{}}
\toprule
关系 & 读课程元数据 & 读课程内容 & 管理课程 \\
\midrule
课程创建者 / 协作教师 & $\checkmark$ & $\checkmark$ & $\checkmark$ \\
已通过审核的学生成员 & $\checkmark$ & $\checkmark$ & $\times$ \\
待审核 / 被拒绝 / 已被移除 & $\checkmark$（仅查看自己的申请状态） & $\times$ & $\times$ \\
公开课程（\textbf{仅学生}） & $\checkmark$ & $\times$ & $\times$ \\
公开课程（教师） & $\times$ & $\times$ & $\times$ \\
无任何关系 & $\times$ & $\times$ & $\times$ \\
\bottomrule
\end{tabularx}
\end{center}
\textbf{两点需要特别说明}：

\begin{enumerate}
\item \textbf{公开课程的元数据只对学生开放。} 学生需要看到公开课程详情才能提交加入申请；教师对本人课程之外无可见性，避免教师误入他人课程。
\item \textbf{「已被移除」优先于「公开」。} 被移除的成员不会因课程转为公开而自动恢复，必须重新申请。
\end{enumerate}

\chapter{系统安装与部署}


\section{快速启动（竞赛演示用）}

\textbf{这是最短路径。} 假设本机已安装 Python 3.10+、Node.js 18+、Neo4j 5.x，完整流程如下：

\begin{Verbatim}[breaklines=true,breakanywhere=true,fontsize=\small,frame=single,framesep=4pt,rulecolor=\color{codeframe}]
# (1) 启动 Neo4j（Neo4j Desktop 点击 Start；或命令行）
cd neo4j-community-5.x && ./bin/neo4j console      # Windows: bin\neo4j.bat console

# (2) 配置后端环境变量
cd backend
cp .env.example .env                                # Windows CMD 见下方对照表
# 用编辑器打开 .env，至少填写：LLM_API_KEY、EMBEDDING_API_KEY、NEO4J_PASSWORD、SECRET_KEY

# (3) 安装后端依赖
python -m venv .venv
.venv\Scripts\activate                              # macOS/Linux: source .venv/bin/activate
pip install -r requirements.txt

# (4) 启动后端（保持终端开启）
python -m uvicorn app.main:app --reload

# (5) 新开终端，启动前端（保持终端开启）
cd frontend-vue
npm install
npm run dev
\end{Verbatim}

\textbf{访问系统：} 浏览器打开 <http://localhost:5173>

\textbf{成功标志：}

\begin{center}
\small
\begin{tabularx}{\linewidth}{@{}XXX@{}}
\toprule
检查点 & 成功标志 & 失败时跳转 \\
\midrule
Neo4j 启动 & 访问 \texttt{http://localhost:7474} 能打开 Neo4j Browser & 6.1 Q1 \\
后端启动 & 终端出现 \texttt{Uvicorn running on http://127.0.0.1:8000} & 6.1 Q2 \\
健康接口 & \texttt{curl http://localhost:8000/health} 返回 \texttt{\{"status":"ok"\}} & 6.1 Q3 \\
前端启动 & 终端出现 \texttt{Local: http://localhost:5173/} & 6.1 Q4 \\
登录页 & 浏览器能打开登录界面 & 6.1 Q5 \\
\bottomrule
\end{tabularx}
\end{center}
\begin{todobox}
\warnmark{} \textbf{前端依赖安装较慢}（含 G6、Element Plus 等）。建议\textbf{提前完成 \texttt{npm install}}，不要在演示现场首次安装。

\end{todobox}


\subsection{跨平台命令对照}

\begin{center}
\small
\begin{tabularx}{\linewidth}{@{}XXXX@{}}
\toprule
操作 & Windows CMD & Windows PowerShell & macOS / Linux \\
\midrule
复制文件 & \texttt{copy .env.example .env} & \texttt{Copy-Item .env.example .env} & \texttt{cp .env.example .env} \\
激活虚拟环境 & \texttt{.venv\textbackslash{}Scripts\textbackslash{}activate} & \texttt{.venv\textbackslash{}Scripts\textbackslash{}Activate.ps1} & \texttt{source .venv/bin/activate} \\
查看端口占用 & \texttt{netstat -ano \textbar{} findstr :8000} & \texttt{Get-NetTCPConnection -LocalPort 8000} & \texttt{lsof -i :8000} \\
结束进程 & \texttt{taskkill /PID <PID> /F} & \texttt{Stop-Process -Id <PID> -Force} & \texttt{kill -9 <PID>} \\
环境变量文件 & 同左 & 同左 & 同左 \\
\bottomrule
\end{tabularx}
\end{center}

\section{运行环境及依赖要求}

\begin{center}
\small
\begin{tabularx}{\linewidth}{@{}XXX@{}}
\toprule
组件 & 版本要求 & 检查命令 \\
\midrule
Python & ≥ 3.10 & \texttt{python --version} \\
Node.js & ≥ 18 & \texttt{node --version} \\
npm & 随 Node.js & \texttt{npm --version} \\
Neo4j & 5.x & Neo4j Desktop 或 \texttt{neo4j --version} \\
\bottomrule
\end{tabularx}
\end{center}
\textbf{硬件与网络要求：}

\begin{center}
\small
\begin{tabularx}{\linewidth}{@{}XX@{}}
\toprule
项 & 要求 \\
\midrule
内存 & 建议 8GB 以上（Neo4j 与两个开发服务器同时运行） \\
磁盘 & ≥ 2GB（含依赖） \\
GPU & \textbf{不需要} \\
网络 & \textbf{必须联网}——知识抽取与向量化均调用外部 API \\
\bottomrule
\end{tabularx}
\end{center}

\section{获取系统与目录结构}

\begin{Verbatim}[breaklines=true,breakanywhere=true,fontsize=\small,frame=single,framesep=4pt,rulecolor=\color{codeframe}]
git clone <仓库地址>          # \verifymark{} 待填：正式提交时替换为实际地址或改为离线交付说明
cd 知识图谱AI分析
\end{Verbatim}

\begin{todobox}
\verifymark{} \textbf{待确认}：若竞赛采用离线交付（U 盘/压缩包），此处的获取方式需相应改写，并说明是否需要一并提供 Neo4j 安装包。

\end{todobox}

\textbf{目录结构：}

\begin{Verbatim}[fontsize=\small,frame=single,framesep=4pt,rulecolor=\color{codeframe}]
知识图谱AI分析/
├── backend/                    后端（FastAPI）
│   ├── app/
│   │   ├── api/                路由层（18 个模块）
│   │   ├── services/           业务逻辑层（24 个模块）
│   │   ├── core/               基础设施（数据库、权限、安全、配置）
│   │   └── utils/              文本处理等工具
│   ├── data/                   运行期数据（SQLite 与上传文件）
│   ├── scripts/                13 个运维/评测脚本
│   ├── test_*.py               端到端验证脚本（自带断言）
│   ├── eval_*.py               评测脚本（结果落 eval_data/）
│   ├── requirements.txt        依赖清单（版本精确锁定）
│   └── .env                    环境变量（由 .env.example 复制）
├── frontend-vue/               前端（Vue 3 + Vite）
│   ├── src/                    页面、组件、路由、状态
│   ├── tests/                  Playwright e2e 测试（11 个 spec）
│   └── vite.config.js          开发服务器与代理配置
└── docs/                       项目文档
\end{Verbatim}


\section{图数据库与外部服务配置}


\subsection{安装并启动 Neo4j}

Neo4j 是\textbf{外部依赖}，不在 pip 依赖内，必须先于后端启动。

\textbf{方式 A：Neo4j Desktop（推荐用于演示环境）}

\begin{enumerate}
\item 官网下载 Neo4j Desktop 并安装
\item 新建 Local DBMS，版本选 5.x
\item 设置密码（后续填入 \texttt{.env}）
\item 点击 \textbf{Start}，确认状态为 Running
\end{enumerate}
\textbf{方式 B：Neo4j Community Server}

\begin{Verbatim}[breaklines=true,breakanywhere=true,fontsize=\small,frame=single,framesep=4pt,rulecolor=\color{codeframe}]
cd neo4j-community-5.x
./bin/neo4j console              # Windows: bin\neo4j.bat console
\end{Verbatim}

\textbf{连接信息：}

\begin{center}
\small
\begin{tabularx}{\linewidth}{@{}XX@{}}
\toprule
项 & 默认值 \\
\midrule
Bolt 地址 & \texttt{bolt://localhost:7687} \\
HTTP 地址 & \texttt{http://localhost:7474} \\
默认用户 & \texttt{neo4j} \\
初始密码 & \texttt{neo4j}（首次登录强制修改） \\
\bottomrule
\end{tabularx}
\end{center}
\begin{notebox}
首次访问 Neo4j Browser 需修改初始密码。\textbf{改完的密码填入 \texttt{.env} 的 \texttt{NEO4J\_PASSWORD}}。

\end{notebox}


\subsection{申请大模型 API Key}

系统需要\textbf{两家}服务商的 Key，缺一不可（原因见下）：

\begin{center}
\small
\begin{tabularx}{\linewidth}{@{}XXXX@{}}
\toprule
用途 & 服务商 & 申请地址 & 说明 \\
\midrule
知识抽取、问答生成 & DeepSeek & platform.deepseek.com &  \\
文本向量化 & SiliconFlow & siliconflow.cn &  \\
\bottomrule
\end{tabularx}
\end{center}
\begin{todobox}
\warnmark{} \textbf{为什么需要两家？} DeepSeek \textbf{没有提供 embedding（向量化）接口}。RAG 检索依赖向量相似度，因此向量化这一步必须使用支持 embedding 的服务（本项目采用 SiliconFlow 的 \texttt{BAAI/bge-m3}）。只配置 DeepSeek 时系统仍可运行，但问答会降级为关键词检索，质量下降。

\end{todobox}


\subsection{配置环境变量}

\begin{Verbatim}[breaklines=true,breakanywhere=true,fontsize=\small,frame=single,framesep=4pt,rulecolor=\color{codeframe}]
cd backend
cp .env.example .env          # Windows CMD: copy .env.example .env
\end{Verbatim}

编辑 \texttt{backend/.env}：

\begin{center}
\small
\begin{tabularx}{\linewidth}{@{}XXX@{}}
\toprule
变量 & 必填 & 说明 \\
\midrule
\texttt{LLM\_API\_KEY} & $\checkmark$ & DeepSeek API Key \\
\texttt{LLM\_API\_BASE} & $\checkmark$ & \texttt{https://api.deepseek.com/v1} \\
\texttt{LLM\_MODEL} & $\checkmark$ & \texttt{deepseek-chat} \\
\texttt{EMBEDDING\_API\_KEY} & $\checkmark$ & SiliconFlow API Key \\
\texttt{EMBEDDING\_API\_BASE} & $\checkmark$ & \texttt{https://api.siliconflow.cn/v1} \\
\texttt{EMBEDDING\_MODEL} & $\checkmark$ & \texttt{BAAI/bge-m3} \\
\texttt{NEO4J\_URI} & $\checkmark$ & \texttt{bolt://localhost:7687} \\
\texttt{NEO4J\_USER} & $\checkmark$ & \texttt{neo4j} \\
\texttt{NEO4J\_PASSWORD} & $\checkmark$ & 2.4.1 中设置的密码 \\
\texttt{SECRET\_KEY} & $\checkmark$ & JWT 签名密钥，见下方生成方法 \\
\texttt{SQLITE\_DB\_PATH} & $\square$ & 默认 \texttt{./data/app.db}（相对 \texttt{backend/}） \\
\texttt{UPLOAD\_DIR} & $\square$ & 默认 \texttt{./data/uploads} \\
\texttt{DEFAULT\_ADMIN\_USERNAME} & $\square$ & 默认 \texttt{sysadmin} \\
\texttt{DEFAULT\_ADMIN\_PASSWORD} & $\square$ & 见 2.6.1 \\
\texttt{DEFAULT\_TEACHER\_PASSWORD} & $\square$ & 见 2.6.1 \\
\bottomrule
\end{tabularx}
\end{center}
\textbf{生成 \texttt{SECRET\_KEY}：}

\begin{Verbatim}[breaklines=true,breakanywhere=true,fontsize=\small,frame=single,framesep=4pt,rulecolor=\color{codeframe}]
python -c "import secrets; print(secrets.token_urlsafe(48))"
\end{Verbatim}

将输出整串粘贴到 \texttt{.env}。\textbf{该密钥泄露等同于任何人都能伪造登录态}，正式部署务必替换模板中的占位值。

\begin{notebox}
\textbf{部署到新环境时的凭据管理建议}：\texttt{.env} 不应提交到代码仓库；演示机器上建议用 \texttt{chmod 600}（Windows 可用文件属性）限制该文件仅运行账号可读。

\end{notebox}


\section{后端与前端部署}


\subsection{后端}

\begin{Verbatim}[breaklines=true,breakanywhere=true,fontsize=\small,frame=single,framesep=4pt,rulecolor=\color{codeframe}]
cd backend
python -m venv .venv
.venv\Scripts\activate            # macOS/Linux: source .venv/bin/activate
pip install -r requirements.txt
\end{Verbatim}

\begin{todobox}
\warnmark{} \textbf{依赖版本已逐个精确锁定，请勿手动升级单个包。} 特别是 \texttt{starlette}：它必须与 \texttt{fastapi 0.115.x} 配套，单独升级会漂移到 1.x 导致后端启动即崩溃。这一点在 \texttt{requirements.txt} 内有注释说明。

\end{todobox}


\subsection{前端}

\begin{Verbatim}[breaklines=true,breakanywhere=true,fontsize=\small,frame=single,framesep=4pt,rulecolor=\color{codeframe}]
cd frontend-vue
npm install
\end{Verbatim}

若安装缓慢或超时：

\begin{Verbatim}[breaklines=true,breakanywhere=true,fontsize=\small,frame=single,framesep=4pt,rulecolor=\color{codeframe}]
npm config set registry https://registry.npmmirror.com
npm install
\end{Verbatim}

前端无需额外配置——\texttt{vite.config.js} 已把 \texttt{/api} 与 \texttt{/health} 代理到 \texttt{http://localhost:8000}。


\section{系统启动与首次登录}


\subsection{启动}

需要\textbf{三个终端}：

\begin{center}
\small
\begin{tabularx}{\linewidth}{@{}XXX@{}}
\toprule
终端 & 命令 & 工作目录 \\
\midrule
(1) Neo4j & 见 2.4.1 & Neo4j 安装目录 \\
(2) 后端 & \texttt{python -m uvicorn app.main:app --reload} & \texttt{backend/} \\
(3) 前端 & \texttt{npm run dev} & \texttt{frontend-vue/} \\
\bottomrule
\end{tabularx}
\end{center}
\textbf{后端首次启动}会自动完成建表、迁移与引导账号创建。\textbf{请留意这段日志}——其中可能包含一次性随机密码。


\subsection{引导账号}

\begin{center}
\small
\begin{tabularx}{\linewidth}{@{}XXXX@{}}
\toprule
账号 & 用户名 & 角色 & 密码来源 \\
\midrule
演示教师账号 & \texttt{admin}（固定） & teacher & \texttt{DEFAULT\_TEACHER\_PASSWORD}，未配置则随机生成 \\
系统管理员账号 & \texttt{sysadmin}（可用 \texttt{DEFAULT\_ADMIN\_USERNAME} 修改） & admin & \texttt{DEFAULT\_ADMIN\_PASSWORD}，未配置则随机生成 \\
\bottomrule
\end{tabularx}
\end{center}
\textbf{两个账号的区别（请勿混淆）：}

\begin{center}
\small
\begin{tabularx}{\linewidth}{@{}XXX@{}}
\toprule
 & 演示教师账号 & 系统管理员账号 \\
\midrule
用途 & 用于演示教学功能（建课、上传、出题） & 用于平台治理（用户、课程、资源、监控） \\
能做什么 & 教学相关全部操作 & 平台级治理操作 \\
不能做什么 & 进入管理工作台 & 参与教学（不属于任何课程） \\
首次登录 & 直接可用 & 使用随机密码时\textbf{强制修改密码}后才可使用 \\
\bottomrule
\end{tabularx}
\end{center}
\textbf{随机密码的获取}：在后端启动日志中搜索 \texttt{演示教师} 或 \texttt{管理员}，日志会打印账号名与密码，\textbf{每个账号只打印一次}。

\begin{notebox}
\textbf{若使用随机密码登录管理员}，系统会强制跳转到「修改密码」页，改完才能进入其它页面。这是设计行为：日志可能被留存或转发，凭据不应长期有效。

\textbf{正式部署请显式配置这两个密码，不要依赖日志传递凭据。} 关于随机密码的日志安全处理，见 7.3 日志及系统监控。

\end{notebox}


\subsection{注册普通账号}

教师与学生请走注册流程：登录页点击「注册」，选择角色（教师 / 学生），自行设置用户名与密码。


\section{部署验收自检表}

\textbf{在新机器上部署完成后，请逐项执行下表并记录结果。} 全部通过才视为部署成功。本表同时可作为竞赛现场环境准备的自检清单。

\begin{center}
\small
\begin{tabularx}{\linewidth}{@{}XXXX@{}}
\toprule
\# & 验收步骤 & 预期结果 & 通过 \\
\midrule
1 & 启动 Neo4j & Neo4j Browser 可打开，Bolt 端口 7687 可连接 & $\square$ \\
2 & 启动后端 & 终端显示 \texttt{Uvicorn running on ...:8000}，无异常堆栈 & $\square$ \\
3 & 检查健康接口 & \texttt{curl http://localhost:8000/health} 返回 \texttt{\{"status":"ok"\}} & $\square$ \\
4 & 启动前端 & 终端显示 \texttt{Local: http://localhost:5173/} & $\square$ \\
5 & 打开登录页 & 浏览器正常加载登录界面 & $\square$ \\
6 & 注册并登录 & 账号创建成功，按角色跳转到正确落点 & $\square$ \\
7 & 创建课程 & 课程创建成功，自动生成 8 位加课码 & $\square$ \\
8 & 上传测试文档 & 解析与抽取均完成，显示知识点数与关系数（均 > 0） & $\square$ \\
9 & 浏览图谱 & 节点与关系正常渲染，点击节点可打开详情 & $\square$ \\
10 & 智能问答 & 返回回答并附带引用来源 & $\square$ \\
11 & 学习路径 & 标记「已掌握」后，推荐结果发生变化 & $\square$ \\
12 & 数据落库检查 & Neo4j 中 \texttt{MATCH (n:KnowledgePoint) RETURN count(n)} 大于 0 & $\square$ \\
\bottomrule
\end{tabularx}
\end{center}
\textbf{第 12 项的说明}：这一步用于确认图谱真的写入了 Neo4j，而不只是 SQLite 里有文档记录。这是新环境部署最容易出错的地方——\textbf{Neo4j 的数据不随代码仓库走}，\texttt{git clone} 只能拿到 SQLite。


\subsection{最小验证数据}

\begin{todobox}
\todomark{} \textbf{【待补：最小验证数据包】}建议随文档提供一份最小测试资料，供部署者快速验证系统可用：

\begin{center}
\small
\begin{tabularx}{\linewidth}{@{}XX@{}}
\toprule
内容 & 要求 \\
\midrule
测试文档 & 一份知识点密集的课程章节（建议 TXT 或 DOCX，避免扫描版 PDF），约 1–2 页 \\
预期结果 & 该文档应抽取出的知识点数量、关系类型分布的\textbf{参考值} \\
存放位置 & 建议置于仓库 \texttt{docs/sample/} 目录 \\
\bottomrule
\end{tabularx}
\end{center}
说明：参考值应来自实际运行结果，用于判断「系统是否正常工作」，不要求完全一致（LLM 抽取存在一定随机性）。

\end{todobox}


\section{赛题指标自检表}

下表将 A10 赛题的技术指标与本文档章节、验证材料对应。\textbf{数值一律标注来源}。

\begin{center}
\small
\begin{tabularx}{\linewidth}{@{}XXXX@{}}
\toprule
赛题指标 & 要求 & 当前状态 & 证据/位置 \\
\midrule
支持文档格式 & ≥ 2 种 & \textbf{4 种}（PDF / DOCX / TXT / MD） & 代码：\texttt{backend/app/services/document\_service.py} \\
知识点数量 & ≥ 20 个（一章内容） & 待实测确认 \todomark{} & 见 5.2；历史文档最大抽取 69 个知识点 \\
关系类型 & ≥ 3 种 & \textbf{4 种}（CONTAINS / PRECEDES / APPLIES\_TO / RELATED\_TO） & 见 4.2 \\
抽取准确率 & ≥ 70\%（人工抽样） & \textbf{已测}：关系主指标 F1 = \textbf{0.8312} & 《知识抽取准确率测试报告》；原始数据 \texttt{backend/eval\_data/eval\_report\_第7章\_树\_v15.json} \\
文档解析+抽取 & ≤ 60 秒 & \textbf{待实测} \todomark{} & 采集方法见 5.4 \\
问答响应 & ≤ 15 秒 & \textbf{待实测} \todomark{} & 采集方法见 5.4 \\
多课程管理 & ≥ 2 门独立课程 & \textbf{待补验证记录} \todomark{} & 见 5.5 \\
普通电脑可运行 & 无需 GPU & 支持 & 见 2.2 \\
\bottomrule
\end{tabularx}
\end{center}

\subsection{关于准确率数据的说明}

上述 \textbf{F1 = 0.8312} 来自已有测试报告，测试条件为：

\begin{center}
\small
\begin{tabularx}{\linewidth}{@{}XX@{}}
\toprule
项 & 值 \\
\midrule
测试文本 & 《Hello 算法》第 7 章「树」（CC BY-NC-SA 4.0），清洗后 15330 字符 \\
分块参数 & chunk 6000 tokens / overlap 400 tokens，并发 4 路 \\
金标准规模 & 49 实体 / 33 关系（人工标注，预测前冻结） \\
Prompt 版本 & v1.5（定稿） \\
评测实现 & \texttt{backend/eval\_accuracy.py}（全仓库唯一评测实现） \\
\bottomrule
\end{tabularx}
\end{center}
\begin{todobox}
\verifymark{} \textbf{两点需在正式报告中明确说明}（属诚实披露，也是赛题「需在文档中说明改进方向」的要求）：

\begin{enumerate}
\item \textbf{该金标准中 PRECEDES 关系为 0 条}（33 条关系由 CONTAINS 25 / RELATED\_TO 5 / APPLIES\_TO 3 构成）。也就是说，\textbf{F1 = 0.8312 覆盖了三种关系类型，但未覆盖前置关系}。而前置关系恰恰是学习路径推荐的基础。建议补充一份包含 PRECEDES 的测试集，否则「三种关系类型均达标」这一表述会不够严谨。
\item \textbf{赛题要求的是「人工抽样验证」}，而当前评测是「对冻结金标准自动计算」。二者并不等价，建议补做一次随机抽样人工判定（如随机抽 30 个知识点），并在报告中同时给出两种口径的结果。
\end{enumerate}
\end{todobox}


\section{生产环境部署与安全建议}

竞赛演示按 2.1 启动即可。若需长期运行或对外提供服务，注意以下几点。


\subsection{服务部署}

\begin{enumerate}
\item \textbf{不要使用 \texttt{--reload}}。它面向开发，会监听文件变化重启，且性能更低。
\item \textbf{前端打包}：\texttt{npm run build} 产出 \texttt{frontend-vue/dist/}，交由 Nginx 等服务托管；此时需把 Nginx 的 \texttt{/api} 反向代理到后端 8000 端口，替代 Vite 的开发代理。
\item \textbf{多进程部署的适用条件}：
\end{enumerate}
\begin{Verbatim}[breaklines=true,breakanywhere=true,fontsize=\small,frame=single,framesep=4pt,rulecolor=\color{codeframe}]
   uvicorn app.main:app --host 0.0.0.0 --port 8000 --workers 4
\end{Verbatim}

\warnmark{} \textbf{不要不加分析地照搬 \texttt{--workers 4}。} 适用性取决于两点：

\begin{itemize}
\item \textbf{SQLite 的写入并发}：SQLite 采用库级写锁，多进程并发写入会互相阻塞甚至报 \texttt{database is locked}。本项目为写入较少的教学场景，若并发写压力上升，应改用「单写进程 + 读副本」或迁移到 PostgreSQL/MySQL；
\item \textbf{Neo4j 连接池}：每个 worker 会建立自己的连接池。worker 数量 × 连接数不应超过 Neo4j 的 \texttt{dbms.connector.bolt.thread\_pool\_max\_size}。
\end{itemize}
在本项目的预期负载下（课堂规模、数十人并发），\textbf{单进程配合异步即可满足}；\texttt{--workers} 主要用于提高并发请求吞吐，需结合实际压测结果确定。

\begin{enumerate}
\item \textbf{收紧 CORS}：当前后端 \texttt{allow\_origins=["*"]} 全放行（便于开发）。对外部署必须改为具体域名。
\item \textbf{启用 HTTPS}：由反向代理终止 TLS。
\item \textbf{显式配置引导密码}，避免依赖启动日志。
\end{enumerate}

\subsection{数据安全}

\begin{center}
\small
\begin{tabularx}{\linewidth}{@{}XX@{}}
\toprule
措施 & 说明 \\
\midrule
\texttt{.env} 权限 & 限制为仅运行账号可读，且\textbf{不纳入版本控制} \\
\texttt{SECRET\_KEY} 轮换 & 泄露后必须更换；更换会使全部已签发 Token 失效，用户需重新登录 \\
日志访问控制 & 见 7.3 \\
定期备份 & 见 7.1 \\
\bottomrule
\end{tabularx}
\end{center}

\subsection{错误信息暴露}

\warnmark{} \textbf{当前实现会把后端异常信息返回给前端。} 这在开发与竞赛演示时有助于定位问题，但\textbf{在生产环境可能向普通用户暴露技术细节}（如文件路径、内部调用栈摘要、外部服务报错原文）。

\textbf{部署建议}：对普通用户，前端应把此类错误统一呈现为友好的提示（如「操作失败，请稍后重试」），详细错误仅写入服务端日志供管理员排查。\textbf{若当前版本尚未做此区分，请在部署说明中如实标注，并作为后续改进项。}


\chapter{系统功能与操作指南}

本章按「操作入口 → 操作步骤 → 预期结果」描述每个功能，并预留界面截图位置。


\section{系统登录与导航}


\subsection{登录}

\textbf{操作入口}：浏览器访问 <http://localhost:5173>

\textbf{操作步骤}：

\begin{enumerate}
\item 在登录页输入用户名与密码
\item 点击「登录」
\end{enumerate}
\textbf{预期结果}：按角色跳转到对应落点——教师与学生进入「课程中心」，管理员进入「管理工作台」。


\subsection{界面导航}

登录后进入带左侧导航的主框架。

\textbf{重要说明}：\textbf{左侧边栏是主要导航入口}。页面顶部的标签栏被刻意隐藏，以避免与侧边栏功能重复。

\begin{center}
\small
\begin{tabularx}{\linewidth}{@{}XX@{}}
\toprule
角色 & 侧边栏菜单 \\
\midrule
教师 & 课程中心 / 数据总览 / 课程管理 / 图谱管理 / 题库管理 / 个人中心 \\
学生 & 课程中心 / 学习总览 / 学习空间（课程文档·图谱浏览·学习路径推荐·收藏夹·做题练习）/ 个人中心 \\
管理员 & 工作台 / 用户管理 / 课程管理 / 资源管理 / 课程治理 / 系统监控 / 审计日志 / 个人中心 \\
\bottomrule
\end{tabularx}
\end{center}
\textbf{全局 AI 助手}：教师端与学生端右下角有常驻的 AI 助手悬浮球，在任意页面均可唤起问答，无需跳转。


\section{教师端操作指南}


\subsection{注册与登录}

\textbf{操作入口}：登录页 → 「注册」

\textbf{操作步骤}：

\begin{enumerate}
\item 点击「注册」进入注册页
\item 填写用户名、密码
\item \textbf{角色选择「教师」}
\item 提交
\end{enumerate}
\textbf{预期结果}：注册成功后自动登录，进入课程中心。


\subsection{创建课程}

\textbf{操作入口}：左侧导航 → \textbf{课程中心} → 「我的课程」

\textbf{操作步骤}：

\begin{enumerate}
\item 点击「新建课程」
\item 填写课程名称、简介等信息
\item 保存
\end{enumerate}
\textbf{预期结果}：

\begin{itemize}
\item 课程卡片出现在「我的课程」列表中
\item 课程自动生成一个 \textbf{8 位加课码}
\end{itemize}
\textbf{关于加课码}：字母表已剔除 \texttt{0/O}、\texttt{1/I/L} 等手写易混字符，便于口头传达给学生。点击卡片上的「邀请」还可生成邀请链接。


\subsection{上传课程文档（核心步骤）}

\textbf{操作入口}：左侧导航 → \textbf{课程管理} → 选择课程 → 「课程文档」

\textbf{操作步骤}：

\begin{enumerate}
\item 选中目标课程
\item 切到「课程文档」页签
\item 点击上传，选择文件
\end{enumerate}
\textbf{支持的格式}：PDF、DOCX、TXT、MD（共 4 种）

\textbf{上传后系统自动执行}：

\begin{Verbatim}[breaklines=true,breakanywhere=true,fontsize=\small,frame=single,framesep=4pt,rulecolor=\color{codeframe}]
解析文档 → 按章节切分与分块 → 调用大模型抽取知识点与关系
        → 写入 Neo4j 图谱 → 回填抽取统计
\end{Verbatim}

\textbf{处理状态}：

\begin{center}
\small
\begin{tabularx}{\linewidth}{@{}XX@{}}
\toprule
状态字段 & 取值过程 \\
\midrule
解析状态 & \texttt{UPLOADED} → \texttt{PARSING} → \texttt{PARSED} / \texttt{FAILED} \\
抽取状态 & \texttt{PENDING} → \texttt{EXTRACTING} → \texttt{COMPLETED} / \texttt{FAILED} \\
\bottomrule
\end{tabularx}
\end{center}
\textbf{预期结果}：处理完成后，文档条目上显示抽取出的\textbf{知识点数量}与\textbf{关系数量}（两者均应大于 0）。

\textbf{异常情况}：

\begin{center}
\small
\begin{tabularx}{\linewidth}{@{}XX@{}}
\toprule
现象 & 原因与处理 \\
\midrule
提示「文件格式不支持」 & 仅支持上述 4 种。注意 \texttt{.doc}（旧版 Word）不支持，需另存为 \texttt{.docx} \\
提示「文档解析结果为空」 & 多为扫描版 PDF（无文本层）。见 6.2 Q1 \\
提示「知识抽取失败」 & 检查 API Key、余额与网络。见 6.5 Q1 \\
\bottomrule
\end{tabularx}
\end{center}
\begin{todobox}
\shotmark{} \textbf{【待补截图 3-1】文档上传与处理结果}

拍摄要点：需体现(1)上传入口 (2)处理中状态 (3)处理完成后的知识点数与关系数。\textbf{这是赛题「至少两种文档格式」的直接证据}，建议同时截取一次 PDF 与一次 DOCX 的上传成功记录。

\end{todobox}

\begin{todobox}
\warnmark{} \textbf{演示提示}：抽取耗时取决于文档长度与 API 响应速度。\textbf{演示前请提前处理好文档}，避免现场长时间等待。建议选用「一章内容」规模的文档——知识点密度高、处理快、讲解也完整。

\end{todobox}


\subsection{查看与编辑知识图谱}

\textbf{操作入口}：左侧导航 → \textbf{图谱管理}


\subsubsection{（1）查看图谱}

\textbf{操作步骤}：

\begin{enumerate}
\item 选择课程与文档
\item 图谱自动加载
\end{enumerate}
\textbf{四种布局}（可切换）：

\begin{center}
\small
\begin{tabularx}{\linewidth}{@{}XX@{}}
\toprule
布局 & 适用场景 \\
\midrule
树图（分层） & 观察 CONTAINS 包含层级结构 \\
环图（同心圆） & 按连接度分布，突出核心概念 \\
个性化（辐射） & 以某个节点为中心查看其邻域 \\
网图（力导向，默认） & 观察整体关联，适合概览 \\
\bottomrule
\end{tabularx}
\end{center}
\textbf{交互操作}：

\begin{center}
\small
\begin{tabularx}{\linewidth}{@{}XX@{}}
\toprule
操作 & 方式 \\
\midrule
缩放 & 鼠标滚轮 \\
平移 & 在空白处拖拽 \\
移动节点 & 拖拽节点 \\
查看详情 & 单击节点，右侧抽屉显示类别、描述与关联关系 \\
定位知识点 & 顶部搜索框输入名称 \\
\bottomrule
\end{tabularx}
\end{center}
\textbf{预期结果}：节点与关系正常渲染，点击节点后详情抽屉正确显示该知识点的信息。


\subsubsection{（2）编辑图谱}

\textbf{操作入口}：图谱管理 → 切换到编辑态

\textbf{可执行的操作}：

\begin{center}
\small
\begin{tabularx}{\linewidth}{@{}XX@{}}
\toprule
操作 & 说明 \\
\midrule
修改知识点 & 改名称、类别、描述 \\
删除知识点 & 删除该节点及其关联关系 \\
新建关系 & 在两个知识点之间建立关系，选择关系类型 \\
\bottomrule
\end{tabularx}
\end{center}
\textbf{四种关系类型}：

\begin{center}
\small
\begin{tabularx}{\linewidth}{@{}XXX@{}}
\toprule
类型 & 中文名 & 含义 \\
\midrule
\texttt{CONTAINS} & 包含 & 上下位、整体-部分、类别-成员 \\
\texttt{PRECEDES} & 前置 & 学习 A 是学习 B 的前提（\textbf{路径推荐依赖此关系}） \\
\texttt{APPLIES\_TO} & 应用于 & A 是方法/算法，用于解决 B 所代表的问题 \\
\texttt{RELATED\_TO} & 相关 & 密切相关但非以上三类 \\
\bottomrule
\end{tabularx}
\end{center}
\begin{todobox}
\warnmark{} \textbf{两点注意}：

\begin{enumerate}
\item \textbf{不要试图创建自定义关系类型。} 系统对关系类型有白名单校验，非法值会回退为「相关关系」，自定义类型不会生效。
\item \textbf{修改知识点描述会使该文档的向量索引失效}，下次问答会重新计算整个文档的向量（涉及外部 API 调用），首次响应会明显变慢。详见 4.4.4。
\end{enumerate}
\end{todobox}

\begin{todobox}
\shotmark{} \textbf{【待补截图 3-2】知识图谱浏览}

拍摄要点：一张能体现\textbf{节点、关系、方向箭头}的图谱全景。建议同时附上节点详情抽屉。

\shotmark{} \textbf{【待补截图 3-3】图谱编辑（修改前后对比）}

拍摄要点：\textbf{这是赛题「教师手动修正图谱」加分项的直接证据}。建议拍成对比图：左「编辑前」、右「编辑后」，体现新增或修改了一个知识点/关系。

\end{todobox}


\subsection{成员管理}

\textbf{操作入口}：左侧导航 → \textbf{课程管理} → 「成员管理」

\textbf{可执行的操作}：

\begin{center}
\small
\begin{tabularx}{\linewidth}{@{}XX@{}}
\toprule
操作 & 说明 \\
\midrule
审核申请 & 查看待审核的学生申请，通过或拒绝 \\
移除成员 & 将已通过成员移出课程 \\
调整角色 & 在「学生」与「协作教师」之间调整 \\
分享课程 & 提供加课码或邀请链接 \\
\bottomrule
\end{tabularx}
\end{center}
\textbf{预期结果}：成员状态变更后，该成员对课程内容的访问权限立即同步。

\begin{notebox}
\textbf{注意}：被移除的成员\textbf{不会}因课程转为公开而自动恢复，必须重新申请。这是刻意的设计（状态优先级）。

\end{notebox}


\subsection{题库管理}

\textbf{操作入口}：左侧导航 → \textbf{题库管理}

\textbf{支持的题型}：单选题、多选题、判断题、填空题、解答题（共 5 种）

\textbf{主要功能}：

\begin{center}
\small
\begin{tabularx}{\linewidth}{@{}XX@{}}
\toprule
功能 & 操作说明 \\
\midrule
手动出题 & 填写题干、选项、正确答案、解析，并关联知识点 \\
自动标注 & 对已有题目点击「自动标注」，系统推荐关联知识点 \\
从文档导入 & 选择一份课程文档，系统尝试从中解析题目 \\
查看统计 & 题量、题型分布、作答正确率、收藏数 \\
\bottomrule
\end{tabularx}
\end{center}
\textbf{关于自动标注}：自动标注的结果\textbf{只填入表单，需点击「保存」才真正写入题库}，并且\textbf{不会覆盖}教师已有的判断。

\textbf{答案可见性}：答案与解析\textbf{仅教师可见}。学生接口经过字段白名单过滤，正确答案只在提交作答后随判分结果返回。


\subsection{学情监测}

\textbf{操作入口}：左侧导航 → \textbf{课程管理} → 「学情监测」

\textbf{操作步骤}：选择课程与文档，查看统计。

\textbf{可查看的内容}：

\begin{center}
\small
\begin{tabularx}{\linewidth}{@{}XX@{}}
\toprule
指标 & 说明 \\
\midrule
已开始学习人数 & 该课程下有学习记录的学生数 \\
平均学习进度 & 全班平均掌握比例 \\
课程知识点总数 & 该文档抽取出的知识点数 \\
学生进度列表 & 按学生列出进度，\textbf{按进度升序排列}（进度落后的排在前面） \\
当前与推荐对比 & 每个学生当前所处知识点与系统推荐的下一步 \\
\bottomrule
\end{tabularx}
\end{center}
\textbf{进度口径}：\texttt{已掌握知识点数 / 知识点总数 × 100\%}


\subsection{作业批改}

\textbf{操作入口}：左侧导航 → \textbf{课程管理} → 「作业批改」，或从题库页的「待批改」卡片进入

\textbf{说明}：客观题（单选/多选/判断）由系统自动判分；主观题（填空/解答）需教师人工评分。


\subsection{文档在线阅读}

\textbf{操作入口}：课程文档列表 → 点击「在线阅读」

\textbf{功能说明}：进入整页阅读器，支持 PDF / DOCX / TXT / MD 四种格式渲染。

\begin{center}
\small
\begin{tabularx}{\linewidth}{@{}XX@{}}
\toprule
区域 & 功能 \\
\midrule
左侧 & 文档目录，可快速跳转章节 \\
右侧 & 本文档知识点面板 \\
顶部 & 返回、文档切换 \\
\bottomrule
\end{tabularx}
\end{center}
\textbf{联动特性（本系统的特色功能）}：

\begin{itemize}
\item 滚动文档时，知识点面板会\textbf{自动高亮}当前阅读位置对应的知识点；
\item 选中正文文本后，可通过浮动菜单\textbf{把选中内容关联到知识点}或加入收藏。
\end{itemize}
\begin{todobox}
\shotmark{} \textbf{【待补截图 3-4】文档在线阅读与知识点联动}

拍摄要点：需同时体现(1)文档正文 (2)右侧知识点面板 (3)当前段落与知识点的对应高亮。\textbf{这是「知识点可追溯到原文」这一特色的直接证据}，也是赛题提交材料第 5 项「图谱截图 + 原始文档片段」所需的素材。

\end{todobox}


\section{学生端操作指南}


\subsection{加入课程}

\textbf{操作入口}：左侧导航 → \textbf{课程中心} → 「加入课程」

\textbf{两种方式}：

\begin{center}
\small
\begin{tabularx}{\linewidth}{@{}XX@{}}
\toprule
方式 & 操作步骤 \\
\midrule
加课码 & 输入教师提供的 8 位加课码，提交申请 \\
邀请链接 & 直接打开教师发来的邀请链接，在落地页点击「申请加入」 \\
\bottomrule
\end{tabularx}
\end{center}
\textbf{预期结果}：提交后状态为\textbf{待审核}，需教师通过后才能访问课程内容。

\begin{notebox}
\textbf{关于加课码输入}：系统会先做归一化处理（去除所有空白并转为大写），因此从聊天软件复制时带上的空格不会导致失败。

\end{notebox}


\subsection{浏览图谱与标记掌握}

\textbf{操作入口}：左侧导航 → \textbf{学习空间} → 「图谱浏览」

\textbf{操作步骤}：

\begin{enumerate}
\item 选择课程与文档
\item 图谱加载后，点击节点查看知识点详情
\item 在详情面板中点击\textbf{「已掌握」}，标记该知识点
\end{enumerate}
\textbf{预期结果}：

\begin{itemize}
\item 被标记的知识点在学习进度中立即体现
\item 该标记会成为\textbf{学习路径推荐的输入}
\end{itemize}
\begin{todobox}
\warnmark{} \textbf{这是使用路径推荐的前置条件。} 没有标记任何已掌握知识点时，系统无法推算你的学习位置。

\end{todobox}


\subsection{学习路径推荐}

\textbf{操作入口}：左侧导航 → \textbf{学习空间} → 「学习路径推荐」

该页面包含两部分功能：


\subsubsection{（1）下一步推荐}

系统基于知识图谱中的前置关系（\texttt{PRECEDES}）自动计算。

\textbf{算法说明}：

\begin{Verbatim}[breaklines=true,breakanywhere=true,fontsize=\small,frame=single,framesep=4pt,rulecolor=\color{codeframe}]
(1) 找出「已掌握知识点所指向、且尚未掌握」的知识点，作为候选
(2) 对每个候选，检查它的全部前置知识是否都已掌握
(3) 仅保留前置全部满足的候选（即：所有前置都必须掌握，而非部分掌握）
(4) 按前置知识数量从多到少排序
(5) 最多返回 10 条
\end{Verbatim}

\textbf{界面展示}：按视觉层级分段——\textbf{已掌握 → 当前 → 推荐 → 未学习}。

\textbf{边界情况}：

\begin{center}
\small
\begin{tabularx}{\linewidth}{@{}XX@{}}
\toprule
情况 & 系统行为 \\
\midrule
尚未标记任何已掌握知识点 & 推荐图谱的\textbf{入门知识点}（没有任何前置关系的节点），最多 10 个 \\
已掌握的知识点未解锁任何新知识点 & 返回该范围内除已掌握外的入门知识点，并提示「未解锁新知识点」 \\
图谱中没有前置关系 & 无候选，退化为上述两种情况之一 \\
\bottomrule
\end{tabularx}
\end{center}
详细的算法约定、方向定义与已知限制见 4.3。


\subsubsection{（2）目标路径}

\textbf{操作步骤}：在「到达目标知识点的学习路径」输入框中填写目标知识点

\textbf{预期结果}：系统生成\textbf{从入门到该目标点}的学习路径，最多返回 5 条，按路径长度升序排列（最短路径优先）。

\begin{notebox}
\textbf{重要说明}：目标路径\textbf{不考虑你已掌握的知识点}，始终返回从入门根节点出发的完整路径。它回答的是「学会这个知识点需要经过哪些步骤」，而不是「我还需要学什么」。

\end{notebox}

\begin{todobox}
\shotmark{} \textbf{【待补截图 3-5】学习路径推荐}

拍摄要点：\textbf{建议拍成前后对比}——(1)标记已掌握之前 (2)标记之后，体现推荐结果随掌握情况变化。这是赛题「已掌握状态变化前后的推荐结果」所需的证据。

\end{todobox}


\subsection{智能问答}

\textbf{操作入口}：

\begin{center}
\small
\begin{tabularx}{\linewidth}{@{}XX@{}}
\toprule
方式 & 位置 \\
\midrule
AI 助手悬浮球（推荐） & 任意页面右下角 \\
智能问答面板 & 访问 \texttt{/student?tab=qa} \\
\bottomrule
\end{tabularx}
\end{center}
\textbf{操作步骤}：

\begin{enumerate}
\item 选择课程与文档（问答是文档级范围的，需先选定）
\item 输入问题
\item 查看回答与\textbf{引用来源}
\end{enumerate}
\textbf{预期结果}：返回基于课程内容的回答，并附上本次回答所依据的课程知识点。

\textbf{技术原理与准确性说明}：见 4.4。请注意该节对问答能力边界的客观描述。

\begin{todobox}
\shotmark{} \textbf{【待补截图 3-6】智能问答与引用来源}

拍摄要点：需同时体现(1)提出的问题 (2)返回的回答 (3)引用来源的知识点列表。\textbf{这是赛题「问答测试集与结果」的配套证据。}

\end{todobox}


\subsection{做题练习与错题本}

\textbf{操作入口}：左侧导航 → \textbf{学习空间} → 「做题练习」

\textbf{功能}：

\begin{center}
\small
\begin{tabularx}{\linewidth}{@{}XX@{}}
\toprule
功能 & 说明 \\
\midrule
选择练习来源 & 路径新知识、随机练习、指定知识点等 \\
作答与判分 & 客观题提交后立即出分；主观题等待教师批改 \\
错题本 & 答错的题自动收录，可针对性重做 \\
\bottomrule
\end{tabularx}
\end{center}
\textbf{判分规则}：

\begin{center}
\small
\begin{tabularx}{\linewidth}{@{}XX@{}}
\toprule
题型 & 规则 \\
\midrule
单选题 & 归一化后比对答案 \\
多选题 & 按\textbf{集合相等}判定，与选项顺序无关；少选或多选均判错 \\
判断题 & 归一化后比对 \\
填空 / 解答 & 由教师人工批改 \\
\bottomrule
\end{tabularx}
\end{center}
\textbf{答案文本的归一化处理}：判分前会统一去除首尾空白、全角转半角、转为小写，因此大小写与全角半角的差异不会影响判分。


\subsection{收藏夹}

\textbf{操作入口}：左侧导航 → \textbf{学习空间} → 「收藏夹」

\textbf{操作方式}：在知识点详情或文档阅读页点击收藏图标，即可收藏知识点或文档。


\subsection{学习总览}

\textbf{操作入口}：左侧导航 → \textbf{学习总览}

这是学生的个人驾驶舱，展示：知识点总数、已掌握数、学习进度、我的收藏数，以及「继续学习」与「学习路径」的快捷入口。


\section{管理员端操作指南}

管理员登录后进入\textbf{管理工作台}，侧边栏包含七个功能模块。


\subsection{工作台}

\textbf{操作入口}：左侧导航 → 「工作台」

\textbf{功能}：平台整体数据概览（用户数、课程数、资源数等）。


\subsection{用户管理}

\textbf{操作入口}：左侧导航 → 「用户管理」

\begin{center}
\small
\begin{tabularx}{\linewidth}{@{}XX@{}}
\toprule
操作 & 说明 \\
\midrule
查看用户列表 & 支持筛选与查询 \\
查看用户详情 & 查看单个账号的基本信息与状态 \\
\textbf{启用 / 禁用账号} & 禁用后该账号无法登录 \\
调整用户角色 & 在教师 / 学生等角色间调整 \\
\textbf{重置密码} & 为指定用户重置密码，用户下次登录需修改 \\
删除账号 & 移除账号 \\
\bottomrule
\end{tabularx}
\end{center}
\begin{notebox}
\textbf{推荐做法}：用户忘记密码时，\textbf{优先使用「重置密码」功能}，而不是直接操作数据库。详见 7.2 账号与密钥管理。

\end{notebox}


\subsection{课程管理}

\textbf{操作入口}：左侧导航 → 「课程管理」

\textbf{说明}：这是\textbf{全平台}的课程视图（区别于教师的「课程中心」，后者只显示本人课程）。

\begin{center}
\small
\begin{tabularx}{\linewidth}{@{}XX@{}}
\toprule
操作 & 说明 \\
\midrule
查看课程列表与详情 & 全平台课程 \\
课程治理 & 对课程执行治理操作 \\
转让课程 & 变更课程归属 \\
删除课程 & 移除课程 \\
\bottomrule
\end{tabularx}
\end{center}

\subsection{资源管理}

\textbf{操作入口}：左侧导航 → 「资源管理」

\begin{center}
\small
\begin{tabularx}{\linewidth}{@{}XX@{}}
\toprule
操作 & 说明 \\
\midrule
文档资源管理 & 查看全平台文档，可执行删除 \\
抽取任务查看 & 查看文档的抽取任务状态 \\
\bottomrule
\end{tabularx}
\end{center}
\begin{notebox}
\textbf{删除文档时的数据一致性}：管理员删除文档走的是\textbf{独立入口}（要求管理员角色），但\textbf{级联清理的顺序与完整性约束与教师入口完全相同}——不会因管理员操作而留下孤立文件或图谱残留。具体清理顺序见 7.1。

\end{notebox}


\subsection{课程治理}

\textbf{操作入口}：左侧导航 → 「课程治理」

\textbf{功能}：对课程执行下架、归档等治理操作。被下架的课程不再对非成员学生暴露元数据。


\subsection{系统监控}

\textbf{操作入口}：左侧导航 → 「系统监控」

\textbf{功能}：展示请求量、响应耗时、错误数等运行指标。

\begin{notebox}
\textbf{指标口径}：这些数据来自\textbf{内存中的实时计数}，不写入数据库、不落日志。这样设计的目的是避免监控本身成为新的性能负担或故障点。\textbf{因此进程重启后计数会清零}，它反映的是本次运行期间的情况，不提供历史趋势。

\end{notebox}


\subsection{审计日志}

\textbf{操作入口}：左侧导航 → 「审计日志」

\textbf{功能}：查看关键操作的留痕记录，用于追溯。


\section{个人中心及账号管理}

\textbf{操作入口}：左侧导航 → 「个人中心」（三种角色共用同一页面）

\begin{center}
\small
\begin{tabularx}{\linewidth}{@{}XX@{}}
\toprule
页签 & 功能 \\
\midrule
基本资料 & 昵称、头像、简介等 \\
密码管理 & 修改密码 \\
注销账号 & 自我注销 \\
\bottomrule
\end{tabularx}
\end{center}
\textbf{首次登录强制改密}：使用引导账号随机密码登录、或被管理员重置过密码的账号，在完成改密前只能停留在「修改密码」页。这是设计行为，不是故障。


\chapter{核心功能原理与使用说明}

本章解释用户\textbf{需要理解}的工作原理。代码级实现细节请参见《项目详细方案》。


\section{文档解析与知识抽取}


\subsection{处理流程}

\begin{Verbatim}[breaklines=true,breakanywhere=true,fontsize=\small,frame=single,framesep=4pt,rulecolor=\color{codeframe}]
上传文档
   ↓
(1) 解析：按格式调用对应解析库提取纯文本
   ↓
(2) 清洗：统一空白字符、去除页眉页脚等噪声
   ↓
(3) 分块：按章节切分，超长章节按段落切分为约 6000 token 的块
   ↓        （相邻块之间保留约 400 token 的重叠）
(4) 抽取：每个块并发提交给大模型，抽取知识点与关系
   ↓
(5) 合并：跨块结果去重、校验、过滤
   ↓
(6) 入图：写入 Neo4j
\end{Verbatim}


\subsection{为什么需要重叠分块}

长文档必须切分后才能交给大模型处理。但切分会产生一个问题：跨块的指代会断裂。例如第 2 块开头的「该算法」指的是第 1 块末尾提到的内容，没有重叠时模型无从判断。

因此相邻块之间保留约 400 token 的重叠区。代价是重叠部分会被重复抽取，由合并环节的去重逻辑兜底。


\subsection{关系类型的判定规则}

系统让大模型按一个\textbf{固定顺序}判断两个知识点之间的关系，先命中先返回：

\begin{center}
\small
\begin{tabularx}{\linewidth}{@{}XXX@{}}
\toprule
顺序 & 判定条件 & 输出关系 \\
\midrule
1 & 上下位 / 整体-部分 / 类别-成员 & \texttt{CONTAINS} \\
2 & 学习 A 是理解 B 的必要前提 & \texttt{PRECEDES} \\
3 & A 是方法/算法，用于解决 B 代表的问题 & \texttt{APPLIES\_TO} \\
4 & 明确密切相关但非上述三类 & \texttt{RELATED\_TO} \\
5 & 证据不足 & \textbf{不输出} \\
\bottomrule
\end{tabularx}
\end{center}
\textbf{特别注意两类容易混淆的情况}：

\begin{itemize}
\item 「栈是一种操作受限的线性表」→ 应判为\textbf{包含关系}（线性表包含栈），而非前置关系；
\item 「本章先介绍线性表，再介绍栈」→ 这只是\textbf{叙述顺序}，不能判为前置关系。
\end{itemize}
因为文本中出现顺序不等于学习依赖顺序，如果把叙述顺序误判为前置关系，学习路径推荐会产生错误结果。

\textbf{设计取向}：系统在提示词中明确要求「宁可少抽取，也不要臆造关系」——即\textbf{优先保证准确率}，允许一定程度的漏抽。


\subsection{合并与校验}

抽取结果在写入图谱前会经过多层校验：

\begin{center}
\small
\begin{tabularx}{\linewidth}{@{}XX@{}}
\toprule
校验项 & 规则 \\
\midrule
实体去重 & 按规范化后的名称去重（全角转半角、折叠空白、纯英文统一大写） \\
关系类型 & 必须属于四种合法类型 \\
自环 & 起点与终点相同的边被丢弃 \\
悬空边 & 两端知识点必须都存在于实体集合中 \\
重复边 & 相同「起点-类型-终点」三元组只保留一条 \\
\bottomrule
\end{tabularx}
\end{center}

\subsection{已识别与未识别的能力边界}

\textbf{已识别的}：

\begin{itemize}
\item 一章规模（约 1.5 万字符）的课程文本，可抽取约 50 个知识点、30 余条关系；
\item 知识点的类别归入四类：概念、定理、公式、方法；
\item 关系类型覆盖四种。
\end{itemize}
\textbf{未识别的（已知限制）}：

\begin{center}
\small
\begin{tabularx}{\linewidth}{@{}XX@{}}
\toprule
限制 & 说明 \\
\midrule
扫描版 PDF & 无文本层，无法解析。需先 OCR \\
图表与公式图像 & 图片内容不参与抽取 \\
跨文档知识融合 & 不提供。同名知识点在不同文档中相互独立 \\
同义词归并 & 当前未启用别名映射表，同一概念的多种写法可能被识别为不同知识点 \\
\bottomrule
\end{tabularx}
\end{center}

\section{知识图谱生成、浏览与编辑}


\subsection{数据组织方式}

\textbf{知识点的唯一标识由三个要素共同确定}：课程 + 文档 + 知识点名称。

这带来的直接结果是：\textbf{同一门课程的两份文档中出现同名知识点时，它们是两个独立的节点}。例如《数据结构》教材和课件中都有「栈」，二者各自拥有独立的描述与关系。

\textbf{为什么这样设计}：修改一份文档的知识点不会波及另一份文档，避免课件修订导致教材图谱被意外改动。


\subsection{与赛题要求的对应}

\begin{center}
\small
\begin{tabularx}{\linewidth}{@{}XX@{}}
\toprule
赛题要求 & 系统实现 \\
\midrule
多课程独立管理 & 按课程隔离，支持同时维护多门课程 \\
不要求跨课程融合 & 课程之间不建立关联，符合赛题要求 \\
文档级隔离 & 同一课程内进一步按文档隔离 \\
\bottomrule
\end{tabularx}
\end{center}

\subsection{图谱的完整性保障}

系统在返回图谱数据前会做\textbf{引用完整性校验}：每条关系的两端必须都能在返回的节点集合中找到，否则该关系被丢弃并记录日志。

\textbf{为什么需要}：前端图渲染库遇到「关系指向不存在的节点」时会直接抛错、整图无法显示。宁可少显示一条边，也不能让整个图谱打不开。


\subsection{编辑操作的影响范围}

\begin{center}
\small
\begin{tabularx}{\linewidth}{@{}XX@{}}
\toprule
编辑动作 & 影响 \\
\midrule
修改知识点名称 / 类别 / 描述 & \textbf{该文档的向量索引整体失效}，下次问答重新计算 \\
删除知识点 & 该节点及其关联关系一并删除 \\
新建 / 删除关系 & 仅影响该关系 \\
\bottomrule
\end{tabularx}
\end{center}
\begin{todobox}
\warnmark{} 第一项影响范围较广，演示与日常使用时需注意，详见 4.4.4。

\end{todobox}


\section{个性化学习路径推荐}


\subsection{关系方向的定义（重要）}

系统中的前置关系有严格的\textbf{方向约定}：

\begin{Verbatim}[breaklines=true,breakanywhere=true,fontsize=\small,frame=single,framesep=4pt,rulecolor=\color{codeframe}]
A -[:PRECEDES]-> B     表示「A 是 B 的前置知识」，即先学 A，再学 B
\end{Verbatim}

由此推出：

\begin{center}
\small
\begin{tabularx}{\linewidth}{@{}XX@{}}
\toprule
概念 & 对应 \\
\midrule
B 的前置知识 & 指向 B 的边（入边） \\
B 解锁的后续知识 & 从 B 出发的边（出边） \\
入门知识点（根节点） & 没有任何入边的节点 \\
\bottomrule
\end{tabularx}
\end{center}
\textbf{在界面上理解}：图谱中箭头从 A 指向 B，表示「学完 A 才能学 B」。


\subsection{下一步推荐的算法}

\begin{Verbatim}[breaklines=true,breakanywhere=true,fontsize=\small,frame=single,framesep=4pt,rulecolor=\color{codeframe}]
输入：学生已掌握的知识点集合 M

(1) 候选生成
   找出所有满足「存在某个已在 M 中的知识点指向它」且「它自身不在 M 中」的节点

(2) 前置饱和检查（核心步骤）
   对每个候选节点，取其全部前置知识集合 P
   仅当 P ⊆ M 时，该候选才被保留
   —— 即：所有前置都必须掌握，而不是掌握其中一部分

(3) 优先级排序
   按 |P|（前置数量）从多到少排序

(4) 截断
   最多返回 10 条
\end{Verbatim}

\textbf{第 (2) 步的意义}：这是「拓扑排序」思想的应用——只有当某个知识点的所有前置都已具备时，学习它才是可行的。把「入度为 0」放宽为「所有前驱已完成」。

\textbf{关于第 (3) 步的说明}：系统按前置数量降序排列，其依据是「前置越多，说明该知识点在知识体系中越靠后、越综合，越应该在当前阶段优先补齐」。\textbf{这是一个业务启发式规则，而非数学上的必然结论。} 其他合理的排序策略（如优先推荐路径最短的知识点）同样成立，本系统的选择已在实现中固定，此处如实说明。


\subsection{目标路径的算法}

\begin{Verbatim}[breaklines=true,breakanywhere=true,fontsize=\small,frame=single,framesep=4pt,rulecolor=\color{codeframe}]
输入：目标知识点 T

沿前置关系从「入门知识点」出发搜索到 T 的路径
  · 起点必须是入门知识点（没有任何前置）
  · 路径深度上限 10 层
  · 按路径长度升序排列，最多返回 5 条（即最短路径优先）
\end{Verbatim}

\textbf{两个重要约定}：

\begin{enumerate}
\item \textbf{起点必须是入门知识点}——排除了「从半路开始」的片段路径，保证推荐路径是完整可学的。
\item \textbf{不考虑学生已掌握的知识点}——无论学生当前掌握到哪里，返回的都是从入门到目标的完整路径。它回答「学会 T 需要经过哪些步骤」，而不是「我还需要学几步」。
\end{enumerate}

\subsection{边界条件与已知限制}

\begin{center}
\small
\begin{tabularx}{\linewidth}{@{}XX@{}}
\toprule
情况 & 系统行为 \\
\midrule
学生未标记任何已掌握知识点 & 推荐入门知识点（无前置的节点），最多 10 个 \\
已掌握但未解锁新知识点 & 返回除已掌握外的入门知识点，提示「未解锁新知识点」 \\
某知识点无前置关系 & 查询前置时，降级返回其「包含」关系的上层概念，标注为「建议先了解的上层概念」 \\
目标知识点无前置路径 & 降级返回该知识点本身 + 其相关概念，提示「该知识点未抽取到先修路径」 \\
目标知识点不存在 & 返回空，前端提示知识点不存在 \\
\bottomrule
\end{tabularx}
\end{center}
\textbf{已知限制（如实说明）}：

\begin{center}
\small
\begin{tabularx}{\linewidth}{@{}XX@{}}
\toprule
限制 & 影响 \\
\midrule
\textbf{无显式环检测} & 系统未对图谱做环检测，防止无限遍历的方式是\textbf{限制路径深度}（前置查询 5 层、目标路径 10 层）。若图谱中存在环，路径不会被无限展开，但可能返回绕行的结果。\textbf{建议教师在建图时避免制造环}（如 A→B→C→A） \\
\textbf{同优先级无确定顺序} & 当前置数量相同时，系统未定义明确的排序规则，显示顺序取决于数据库返回顺序。\textbf{相同优先级的推荐结果之间没有稳定顺序} \\
\textbf{前置查询可能重复} & 若某个前置知识点通过多条路径到达目标，查询结果中可能出现重复条目 \\
\bottomrule
\end{tabularx}
\end{center}
\begin{notebox}
上述限制属于当前版本的客观情况，建议在《项目详细方案》中作为改进方向记录。

\end{notebox}


\subsection{推荐示例}

\begin{quote}
\small
\todomark{} \textbf{【待补：推荐示例】}建议用\textbf{真实课程中的实际知识点}绘制一张示例图，展示：

\begin{Verbatim}[breaklines=true,breakanywhere=true,fontsize=\small,frame=single,framesep=4pt,rulecolor=\color{codeframe}]
已掌握：〔知识点 A〕、〔知识点 B〕
   ↓ 箭头标注关系类型与方向
系统推荐：〔知识点 C〕   推荐理由：其全部前置（A、B）均已掌握，前置数 2
尚未解锁：〔知识点 D〕   原因：其前置包含 C，C 尚未掌握
\end{Verbatim}

要求：注明每条箭头的\textbf{关系方向}、\textbf{掌握条件}与\textbf{推荐依据}，使评委能据此核对算法说明与实际行为是否一致。

\warnmark{} 必须使用系统中真实的课程知识点，不要用虚构示例。

\end{quote}


\section{RAG 智能问答}


\subsection{检索链路}

\begin{Verbatim}[breaklines=true,breakanywhere=true,fontsize=\small,frame=single,framesep=4pt,rulecolor=\color{codeframe}]
学生提问
   ↓
(1) 问题向量化（调用 bge-m3）
   ↓
(2) 取回该文档的全部知识点向量（必要时先构建/更新索引）
   ↓
(3) 计算问题向量与每个知识点向量的余弦相似度，取最相关的若干条
   ↓
(4) 按知识点标识回查图谱，取回知识点的名称、类别、描述与关系
   ↓
(5) 与问题一起提交给大模型生成回答
   ↓
(6) 返回「回答 + 本次使用的知识点引用来源」
\end{Verbatim}


\subsection{关于回答准确性的客观说明}

\textbf{系统通过提示词约束模型优先依据检索到的课程内容作答}，具体约束包括：

\begin{enumerate}
\item 优先使用提供的知识内容回答问题；
\item 如果知识库中没有相关信息，应如实告知，不要编造；
\item 回答要准确、简洁、易于理解；
\item 如合适，可推荐相关知识点供进一步学习。
\end{enumerate}
\textbf{需要明确的技术边界}：

\begin{center}
\small
\begin{tabularx}{\linewidth}{@{}XX@{}}
\toprule
说明 & 内容 \\
\midrule
引用来源的含义 & 引用来源是\textbf{本次实际提供给模型的上下文}，即「回答所依据的材料」。它证明检索到了哪些内容，\textbf{但不等于回答必然正确} \\
无强制校验 & 系统\textbf{未对模型的输出做事实性校验}。即使有检索内容，大模型仍可能生成检索内容未直接支持的信息 \\
检索质量决定上限 & 若检索到的知识点与问题不相关，回答质量会下降 \\
检索不足时的行为 & 提示词要求模型如实告知知识不足，但这是\textbf{提示词层面的软约束}，不是硬性拦截 \\
\bottomrule
\end{tabularx}
\end{center}
\begin{notebox}
\textbf{结论}：系统通过「先检索课程内容、再生成回答、并给出引用来源」的方式降低了无依据生成的风险，提高了回答与课程内容的一致性。\textbf{回答的准确性应通过问答测试集进行验证}，结果见 5.2 与《问答测试集》。

\end{notebox}


\subsection{检索的降级机制}

为保障弱网或未配置向量服务时功能仍可用，系统采用三级降级：

\begin{center}
\small
\begin{tabularx}{\linewidth}{@{}XXX@{}}
\toprule
级别 & 触发条件 & 行为 \\
\midrule
一级：向量检索 & 正常情况 & 语义相似度检索 \\
二级：关键词检索 & 向量检索失败或未配置向量服务 & 按关键词匹配知识点 \\
三级：同文档补充 & 检索结果不足 & 用同文档其他知识点补足上下文 \\
\bottomrule
\end{tabularx}
\end{center}
\begin{todobox}
\warnmark{} \textbf{降级是静默的}。当向量服务不可用时，系统不会报错，而是自动退回关键词检索——表现为「回答质量下降但不报错」。排查时请查看后端日志中的警告信息，或确认 \texttt{.env} 中的 \texttt{EMBEDDING\_API\_KEY} 是否有效。详见 6.5 Q3。

\end{todobox}


\subsection{索引的构建与失效}

\textbf{构建时机}：系统采用\textbf{懒构建}——首次问答时若发现索引缺失，才现场生成该文档全部知识点的向量。

\textbf{重建时机}：每次问答前会比对图谱中当前的知识点与已存储的向量是否一致，不一致则整体重建。

\textbf{失效粒度}：\textbf{整个文档}。教师修改任意一个知识点的描述，会导致该文档全部向量被删除，下次问答时全量重算。

\textbf{为什么要整份重建而不是只更新一条}：知识点重新抽取时，其内部标识会整体更换。旧向量会指向已不存在的节点，而数据库表中的记录「非空」会让系统误以为索引有效、永不重建。因此系统采用「整份比对、整份重建」的策略，用一定的性能代价换取索引不会失效后无声无息。

\textbf{对使用者的影响}：

\begin{center}
\small
\begin{tabularx}{\linewidth}{@{}XX@{}}
\toprule
场景 & 表现 \\
\midrule
首次对该文档提问 & 明显较慢（需现场生成向量） \\
教师编辑过该文档的知识点后首次提问 & 明显较慢（全量重建） \\
后续提问 & 正常速度 \\
\bottomrule
\end{tabularx}
\end{center}
\begin{todobox}
\warnmark{} \textbf{演示提示}：如果演示流程中包含「先修改图谱、再当场提问」，请在修改后\textbf{先提一个问题把索引热起来}，否则那一问的耗时会明显偏长。

\end{todobox}


\subsection{检索方式的性能特征}

系统采用\textbf{全文比对}的方式计算相似度，未使用近似最近邻（ANN）索引。

\begin{center}
\small
\begin{tabularx}{\linewidth}{@{}XX@{}}
\toprule
项 & 说明 \\
\midrule
复杂度 & 与知识点数量成正比 \\
适用规模 & 单文档数百个知识点以内，响应时间为毫秒级 \\
不适用 & 跨全部课程的大规模检索 \\
\bottomrule
\end{tabularx}
\end{center}
\textbf{为什么不用 ANN}：当前文档规模（数十至数百个知识点）远未达到需要索引的程度，全文比对的结果是\textbf{精确解}而非近似解，反而更准确。相关对比实验见《论文算法复现实验》。


\section{题库、练习及学情分析}


\subsection{判分规则}

\begin{center}
\small
\begin{tabularx}{\linewidth}{@{}XX@{}}
\toprule
题型 & 判分方式 \\
\midrule
单选题 & 文本归一化后比对 \\
多选题 & 集合相等判定（顺序无关） \\
判断题 & 归一化后比对 \\
填空题 & 逐空比对 \\
解答题 & 教师人工批改 \\
\bottomrule
\end{tabularx}
\end{center}
\textbf{归一化规则}：判分前统一去除首尾空白、全角字符转半角、转为小写。因此「Ａ」与「a」等价。


\subsection{答案的可见性控制}

学生端接口采用\textbf{字段白名单}方式构造响应，\texttt{answer}（答案）与 \texttt{analysis}（解析）字段\textbf{不会下发}给学生。正确答案仅在该题提交作答后随判分结果返回。


\subsection{题目删除的策略}

\begin{center}
\small
\begin{tabularx}{\linewidth}{@{}XX@{}}
\toprule
情况 & 行为 \\
\midrule
题目已有作答记录 & \textbf{软删除}：移出出题池，但不删除记录，保护学生作答历史 \\
题目无任何作答记录 & 物理删除 \\
文档被删除时 & 该文档的题目及其作答记录一并清理；\textbf{课程通用题（不属于任何文档）保留} \\
\bottomrule
\end{tabularx}
\end{center}

\subsection{学情分析的进度口径}

\begin{Verbatim}[breaklines=true,breakanywhere=true,fontsize=\small,frame=single,framesep=4pt,rulecolor=\color{codeframe}]
学习进度 = 已掌握知识点数 / 该文档知识点总数 × 100%
\end{Verbatim}

进度\textbf{按文档计算}。切换文档后，进度会相应变化。


\chapter{系统测试与赛题指标验证}

\begin{notebox}
\textbf{本章说明}：所有数值均标注来源。已有实测数据的项目直接引用对应报告；尚无记录的项目留空并注明采集方法，\textbf{不以预期值代替实测值}。

\end{notebox}


\section{测试环境与测试数据}


\subsection{测试环境}

\begin{todobox}
\todomark{} \textbf{【待补：测试环境记录】}建议按下表实测填写，供评委判断结果的可复现性：

\begin{center}
\small
\begin{tabularx}{\linewidth}{@{}XX@{}}
\toprule
项 & 实际配置 \\
\midrule
操作系统 & 待填 \\
CPU / 内存 & 待填 \\
Python / Node.js 版本 & 待填 \\
Neo4j 版本 & 待填 \\
网络环境 & 待填（是否专线、是否有代理） \\
测试时间 & 待填 \\
\bottomrule
\end{tabularx}
\end{center}
\end{todobox}


\subsection{测试数据}

\begin{center}
\small
\begin{tabularx}{\linewidth}{@{}XXX@{}}
\toprule
数据集 & 内容 & 来源 \\
\midrule
知识抽取测试 & 《Hello 算法》第 7 章「树」（CC BY-NC-SA 4.0），清洗后 15330 字符 & \texttt{backend/eval\_data/第7章\_树.txt} \\
抽取金标准 & 49 实体 / 33 关系，人工标注、预测前冻结 & \texttt{backend/eval\_data/gold\_第7章\_树.json} \\
问答测试集 & 21 题，覆盖 6 类问题 & 《问答测试集》 \\
\bottomrule
\end{tabularx}
\end{center}
\begin{notebox}
\textbf{关于教材版权}：赛题明确要求课程资料不得侵犯版权，建议使用开源教材或自编示例。本条所用的《Hello 算法》采用 CC BY-NC-SA 4.0 协议，属开源教材。

\end{notebox}


\section{核心功能测试}

\begin{todobox}
\todomark{} \textbf{【待补：核心功能测试结果】}建议按下表逐项实测并记录。\textbf{「实际结果」一列必须填写真实观察到的现象，不要照抄预期结果。}

\begin{center}
\footnotesize
\begin{tabularx}{\linewidth}{@{}XXXXXX@{}}
\toprule
\# & 测试项 & 预期结果 & 实际结果 & 是否通过 & 备注 \\
\midrule
1 & PDF 文档上传与解析 & 解析成功，抽取知识点 > 0 & 待填 & $\square$ &  \\
2 & DOCX 文档上传与解析 & 解析成功，抽取知识点 > 0 & 待填 & $\square$ &  \\
3 & TXT / MD 文档上传 & 解析成功 & 待填 & $\square$ &  \\
4 & 知识图谱渲染 & 节点与关系正常显示 & 待填 & $\square$ &  \\
5 & 图谱交互 & 缩放、拖拽、点击详情均正常 & 待填 & $\square$ &  \\
6 & 图谱编辑 & 增删节点/关系生效并持久化 & 待填 & $\square$ &  \\
7 & 知识点详情展示 & 定义、类别、关系正确 & 待填 & $\square$ &  \\
8 & 智能问答 & 返回回答与引用来源 & 待填 & $\square$ &  \\
9 & 学习路径推荐 & 标记掌握后推荐结果变化 & 待填 & $\square$ &  \\
10 & 目标路径生成 & 返回完整路径 & 待填 & $\square$ &  \\
11 & 文档在线阅读 & 四种格式渲染正常 & 待填 & $\square$ &  \\
12 & 题库与判分 & 客观题自动判分正确 & 待填 & $\square$ &  \\
13 & 成员审核流程 & 申请-审核-权限生效 & 待填 & $\square$ &  \\
14 & 权限隔离 & 非成员无法访问课程内容 & 待填 & $\square$ &  \\
\bottomrule
\end{tabularx}
\end{center}
\end{todobox}


\subsection{问答效果验证}

《问答测试集》已包含 21 道测试题，测试条件如下：

\begin{center}
\small
\begin{tabularx}{\linewidth}{@{}XX@{}}
\toprule
项 & 值 \\
\midrule
测试课程 & course\_id=5《第6章 面向对象程序设计》 \\
图谱规模 & 69 个知识点、118 条关系（CONTAINS 45 / RELATED\_TO 35 / APPLIES\_TO 32 / PRECEDES 6） \\
向量模型 & SiliconFlow \texttt{BAAI/bge-m3}（1024 维） \\
生成模型 & DeepSeek \texttt{deepseek-chat}（temperature=0.3） \\
检索参数 & top\_k=5 \\
\bottomrule
\end{tabularx}
\end{center}
\begin{todobox}
\todomark{} \textbf{【待补】}：请将《问答测试集》中的\textbf{逐题结果汇总}到本文档，或在此处直接引用其结论章节。建议呈现形式：

\begin{center}
\footnotesize
\begin{tabularx}{\linewidth}{@{}XXXXX@{}}
\toprule
题号 & 问题 & 检索命中的知识点 & 系统回答（摘要） & 人工评价 \\
\midrule
1 & 待填 & 待填 & 待填 & 待填 \\
\bottomrule
\end{tabularx}
\end{center}
人工评价建议采用「完全正确 / 部分正确 / 错误 / 诚实告知知识不足」四档。

\end{todobox}


\section{知识抽取准确率测试}

\textbf{本节结论引自《知识抽取准确率测试报告》，完整方法与过程请参见该报告。}


\subsection{测试条件}

\begin{center}
\small
\begin{tabularx}{\linewidth}{@{}XX@{}}
\toprule
项 & 值 \\
\midrule
测试文本 & 《Hello 算法》第 7 章「树」，15330 字符 \\
分块参数 & chunk 6000 tokens / overlap 400 tokens，并发 4 路 \\
金标准 & 49 实体 / 33 关系，人工标注并在预测前冻结（含 SHA256 记录） \\
Prompt 版本 & v1.5（定稿） \\
评测实现 & \texttt{backend/eval\_accuracy.py}（全仓库唯一评测实现） \\
\bottomrule
\end{tabularx}
\end{center}

\subsection{测试结果}

\begin{center}
\small
\begin{tabularx}{\linewidth}{@{}XXX@{}}
\toprule
指标 & 结果 & 说明 \\
\midrule
关系主指标 F1 & \textbf{0.8312} & 赛题申报口径 \\
实体类别准确率 & 见测试报告 & 在「已识别实体」的交集上评价 \\
\bottomrule
\end{tabularx}
\end{center}
\textbf{关系主指标的统计口径}：

\begin{itemize}
\item Precision = 正确预测数 / 预测总数（金标准全部关系为分母基准）
\item Recall = 正确预测数 / 核心关系数（核心关系 = 原文明确、无争议的关系）
\item 模型的输出必须落在金标准内才算正确；方向错误、类型错误、臆造均全额扣分
\end{itemize}
\textbf{辅助指标}（详见测试报告）：按关系类型的 P/R/F1、方向错误率、错误归因五分类（臆造 / 方向颠倒 / 类型混淆 / 同义拆分 / 漏抽）。

\begin{notebox}
\textbf{说明}：模型输出携带的 \texttt{confidence}（自评置信度）\textbf{不参与任何 P/R/F1 计算}，仅用于阈值过滤。

\end{notebox}


\subsection{与赛题指标的对照}

\begin{center}
\small
\begin{tabularx}{\linewidth}{@{}XXX@{}}
\toprule
赛题要求 & 结果 & 是否达标 \\
\midrule
准确率 ≥ 70\% & F1 = 0.8312 & $\checkmark$ \\
\bottomrule
\end{tabularx}
\end{center}
\begin{todobox}
\verifymark{} \textbf{需在正式报告中补充说明的两点}（见 2.8）：

\begin{enumerate}
\item 该金标准的 33 条关系中 \textbf{PRECEDES 为 0 条}，故当前 F1 未覆盖前置关系；
\item 赛题要求「人工抽样验证」，当前为「对冻结金标准的自动评测」，建议补做人工抽样。
\end{enumerate}
\end{todobox}


\section{文档处理与问答性能测试}

\begin{quote}
\small
\todomark{} \textbf{【待补：性能测试记录】}

\textbf{当前状态}：仓库中有性能测试脚本 \texttt{backend/test\_perf\_timing.py}（输出解析耗时、分块数、抽取耗时与总耗时），但\textbf{尚未留存正式的性能测试记录}。

\textbf{采集方法}：

\begin{Verbatim}[breaklines=true,breakanywhere=true,fontsize=\small,frame=single,framesep=4pt,rulecolor=\color{codeframe}]
cd backend
python test_perf_timing.py
\end{Verbatim}

请在与演示相同的机器、相同的测试文档上运行至少 3 次，记录结果如下：


\textbf{文档解析与抽取耗时（赛题要求 ≤ 60 秒）}\par

\begin{center}
\footnotesize
\begin{tabularx}{\linewidth}{@{}XXXXXXX@{}}
\toprule
轮次 & 文档字符数 & 分块数 & 解析耗时 & 抽取耗时 & 总耗时 & 是否达标 \\
\midrule
第 1 次 & 待填 & 待填 & 待填 & 待填 & 待填 & $\square$ \\
第 2 次 & 待填 & 待填 & 待填 & 待填 & 待填 & $\square$ \\
第 3 次 & 待填 & 待填 & 待填 & 待填 & 待填 & $\square$ \\
\bottomrule
\end{tabularx}
\end{center}

\textbf{智能问答响应耗时（赛题要求 ≤ 15 秒）}\par

\begin{center}
\footnotesize
\begin{tabularx}{\linewidth}{@{}XXXXXX@{}}
\toprule
场景 & 测试题数 & 平均耗时 & P95 耗时 & 最大耗时 & 是否达标 \\
\midrule
索引已就绪 & 待填 & 待填 & 待填 & 待填 & $\square$ \\
首次提问（含索引构建） & 待填 & 待填 & 待填 & 待填 & $\square$ \\
\bottomrule
\end{tabularx}
\end{center}
\begin{notebox}
\textbf{建议分两种场景分别记录。} 首次提问需要现场构建向量索引，耗时显著高于常态。分开记录既更真实，也能说明系统在稳态下的性能。

\end{notebox}

\textbf{注意}：LLM 与向量服务的响应速度受网络与负载影响，存在轮次波动。建议注明测试时间与网络环境。

\end{quote}


\section{多课程隔离与权限测试}

\begin{todobox}
\todomark{} \textbf{【待补：多课程与权限验证记录】}

\textbf{多课程独立管理}（赛题要求 ≥ 2 门独立课程）

\begin{center}
\small
\begin{tabularx}{\linewidth}{@{}XXX@{}}
\toprule
验证项 & 课程 A & 课程 B \\
\midrule
课程名称 & 待填 & 待填 \\
文档数 & 待填 & 待填 \\
知识点数 & 待填 & 待填 \\
关系数 & 待填 & 待填 \\
数据是否互相隔离 & 待验证 & 待验证 \\
\bottomrule
\end{tabularx}
\end{center}
\textbf{验证方法}：在两门课程中分别上传文档并抽取，确认：

\begin{enumerate}
\item 两门课程的知识图谱互不干扰；
\item 在课程 A 上下文中查询不到课程 B 的知识点；
\item 删除课程 A 的文档不影响课程 B 的数据。
\end{enumerate}
\textbf{权限测试}：仓库中已有 \texttt{backend/test\_permission\_matrix.py}（权限矩阵验证脚本），建议运行并留存输出。

\begin{center}
\footnotesize
\begin{tabularx}{\linewidth}{@{}XXXXX@{}}
\toprule
角色关系 & 读元数据 & 读内容 & 管理 & 实际结果 \\
\midrule
课程创建者 & $\checkmark$ & $\checkmark$ & $\checkmark$ & 待填 \\
协作教师 & $\checkmark$ & $\checkmark$ & $\checkmark$ & 待填 \\
已通过学生 & $\checkmark$ & $\checkmark$ & $\times$ & 待填 \\
待审核 & $\checkmark$ & $\times$ & $\times$ & 待填 \\
公开课程（学生） & $\checkmark$ & $\times$ & $\times$ & 待填 \\
无关系 & $\times$ & $\times$ & $\times$ & 待填 \\
\bottomrule
\end{tabularx}
\end{center}
\end{todobox}


\section{测试结果与问题处理记录}

\begin{todobox}
\todomark{} \textbf{【待补：问题处理记录】}

建议记录测试过程中发现的\textbf{真实问题}及其处理结果。如实记录问题并说明处理方式，比只记录成功结果更能体现工程完整性。

\begin{center}
\footnotesize
\begin{tabularx}{\linewidth}{@{}XXXXX@{}}
\toprule
\# & 问题现象 & 原因分析 & 处理方式 & 状态 \\
\midrule
1 & 待填 & 待填 & 待填 & 待填 \\
\bottomrule
\end{tabularx}
\end{center}
仓库中已有 \texttt{docs/问题修复记录.md}，可直接引用其中的典型问题作为素材。

\end{todobox}


\chapter{常见问题与故障排查}

本章按\textbf{使用者角色}分类。每条问题采用统一格式：

\begin{notebox}
\textbf{问题现象} → \textbf{可能原因} → \textbf{解决步骤} → \textbf{仍无法解决时}

\end{notebox}


\section{安装部署问题}


\subsection{Q1：Neo4j 无法启动 / 连接不上}

\textbf{问题现象}：后端启动时报连接错误，或 Neo4j Desktop 状态异常。

\textbf{可能原因}：Neo4j 未启动；密码错误；端口被占用。

\textbf{解决步骤}：

\begin{enumerate}
\item 确认 Neo4j 已启动（Neo4j Desktop 显示 Running）
\item 访问 \texttt{http://localhost:7474}，确认能打开 Neo4j Browser
\item 核对 \texttt{.env} 中的 \texttt{NEO4J\_PASSWORD} 与 Neo4j 实际密码是否一致
\item 检查 7687 端口是否被占用（命令见 2.1 跨平台对照表）
\end{enumerate}
\textbf{仍无法解决时}：查看 Neo4j 的 \texttt{logs/} 目录，看具体报错。常见的是端口冲突或数据目录损坏（后者可尝试新建一个 DBMS）。


\subsection{Q2：后端启动崩溃，报 starlette 相关错误}

\textbf{问题现象}：\texttt{python -m uvicorn app.main:app --reload} 立即退出并打印异常。

\textbf{可能原因}：依赖版本不匹配。\texttt{starlette} 被升级到了与 \texttt{fastapi 0.115.x} 不兼容的版本。

\textbf{解决步骤}：

\begin{Verbatim}[breaklines=true,breakanywhere=true,fontsize=\small,frame=single,framesep=4pt,rulecolor=\color{codeframe}]
cd backend
pip install -r requirements.txt --force-reinstall
\end{Verbatim}

\textbf{仍无法解决时}：确认是否曾执行过 \texttt{pip install -U starlette} 之类的单独升级。重建虚拟环境是最可靠的方案：

\begin{Verbatim}[breaklines=true,breakanywhere=true,fontsize=\small,frame=single,framesep=4pt,rulecolor=\color{codeframe}]
rm -rf .venv          # Windows: rmdir /s .venv
python -m venv .venv
.venv\Scripts\activate
pip install -r requirements.txt
\end{Verbatim}


\subsection{Q3：健康接口返回异常或无法访问}

\textbf{问题现象}：\texttt{curl http://localhost:8000/health} 无响应或返回错误。

\textbf{可能原因}：后端未启动；端口不是 8000。

\textbf{解决步骤}：

\begin{enumerate}
\item 确认后端终端仍在运行且无报错
\item 确认终端输出的端口号（若改过端口，需同步修改 \texttt{vite.config.js} 的代理目标）
\item 重新执行 \texttt{curl http://localhost:8000/health}，期望返回 \texttt{\{"status":"ok"\}}
\end{enumerate}
\textbf{仍无法解决时}：查看后端终端完整输出，定位启动阶段的异常。


\subsection{Q4：前端启动失败或依赖安装报错}

\textbf{问题现象}：\texttt{npm run dev} 报错，或 \texttt{npm install} 失败。

\textbf{可能原因}：Node.js 版本过低；npm 源网络问题。

\textbf{解决步骤}：

\begin{enumerate}
\item 确认 \texttt{node --version} ≥ 18
\item 切换镜像源后重装：
\end{enumerate}
\begin{Verbatim}[breaklines=true,breakanywhere=true,fontsize=\small,frame=single,framesep=4pt,rulecolor=\color{codeframe}]
   npm config set registry https://registry.npmmirror.com
   rm -rf node_modules package-lock.json
   npm install
\end{Verbatim}

\textbf{仍无法解决时}：删除 \texttt{node\_modules} 后重试；若仍失败，检查是否有代理软件拦截 npm 请求。


\subsection{Q5：页面打开是空白}

\textbf{问题现象}：浏览器访问 5173 端口，页面空白。

\textbf{可能原因}：前端未启动；或后端未启动导致接口请求全部失败。

\textbf{解决步骤}：

\begin{enumerate}
\item 按 F12 打开控制台，查看报错
\item 确认后端已启动（见 Q3）
\item 确认前端终端显示 \texttt{Local: http://localhost:5173/}
\end{enumerate}
\textbf{仍无法解决时}：查看控制台的网络请求，确认 \texttt{/api} 请求是否返回 404 或连接被拒。


\subsection{Q6：端口被占用}

\textbf{问题现象}：启动时报 \texttt{address already in use}。

\textbf{解决步骤}：按 2.1 的跨平台命令对照表查找并结束占用进程，或改用其他端口启动。

\textbf{仍无法解决时}：重启机器是最直接的办法。


\section{教师端操作问题}


\subsection{Q1：上传后知识点为 0，或提示「文档解析结果为空」}

\textbf{问题现象}：文档上传成功，但抽取出的知识点数为 0。

\textbf{可能原因}：\textbf{扫描版 PDF 没有文本层}——它本质是图片，无法提取文字。

\textbf{解决步骤}：

\begin{enumerate}
\item 用 PDF 阅读器尝试选中正文文字
\item 如果无法选中，说明是扫描件
\item 换用电子版文档，或先用 OCR 工具转换为可复制文本的版本
\end{enumerate}
\textbf{仍无法解决时}：换一份 DOCX 或 TXT 文档测试，以排除是否格式特有问题。


\subsection{Q2：提示「文件格式不支持」}

\textbf{问题现象}：上传时被拒绝。

\textbf{可能原因}：文件扩展名不在支持范围内，或实际格式与扩展名不符。

\textbf{解决步骤}：

\begin{enumerate}
\item 确认扩展名为 \texttt{.pdf} / \texttt{.docx} / \texttt{.txt} / \texttt{.md} 之一
\item 若是 \texttt{.doc}（旧版 Word），用 Word 另存为 \texttt{.docx}
\item 若文件从其他格式改名而来（如把 \texttt{.doc} 改成 \texttt{.docx}），需用 Word 真正转换
\end{enumerate}
\textbf{仍无法解决时}：换一份确认无误的文档测试。


\subsection{Q3：抽取失败，提示「知识抽取失败」}

\textbf{问题现象}：解析成功但抽取阶段报错。

\textbf{可能原因}：API Key 无效；账号余额不足；网络不通。

\textbf{解决步骤}：

\begin{enumerate}
\item 检查 \texttt{.env} 中的 \texttt{LLM\_API\_KEY} 是否有效
\item 登录 DeepSeek 平台确认账户余额
\item 检查网络能否访问 \texttt{LLM\_API\_BASE}
\item 查看后端终端的具体报错信息
\end{enumerate}
\textbf{仍无法解决时}：见 6.5 Q1。


\subsection{Q4：学生无法加入课程}

\textbf{问题现象}：学生输入加课码后提示无效，或申请一直处于待审核。

\textbf{可能原因}：加课码输入错误；申请尚未被审核。

\textbf{解决步骤}：

\begin{enumerate}
\item 核对加课码（系统已剔除 \texttt{0/O}、\texttt{1/I/L} 等易混字符，不会出现这些字符）
\item 学生提交后，教师需在「成员管理」中\textbf{通过}其申请
\item 检查学生是否曾被移除（被移除的成员需重新申请，不会自动恢复）
\end{enumerate}
\textbf{仍无法解决时}：改用邀请链接的方式让学生加入。


\subsection{Q5：图谱修改失败}

\textbf{问题现象}：编辑知识图谱时保存失败。

\textbf{可能原因}：必填字段缺失；存在同名知识点冲突。

\textbf{解决步骤}：

\begin{enumerate}
\item 确认知识点名称等必填项已填写
\item 若提示重名，检查该文档下是否已存在同名知识点
\end{enumerate}
\textbf{仍无法解决时}：刷新页面后重试，并查看浏览器控制台的报错信息。


\subsection{Q6：想重新处理一份文档}

\textbf{操作建议}：当前版本不支持「重新抽取」按钮。如需重做，请\textbf{删除该文档后重新上传}。

\begin{todobox}
\warnmark{} \textbf{注意}：删除文档会级联清理该文档的知识图谱、向量、学习记录、收藏与该文档的题目。\textbf{学生的相关学习记录会一并删除}，操作前请确认。

\end{todobox}

\textbf{仍无法解决时}：如仅需修正个别知识点，可在图谱编辑态直接修改，无需重新上传。


\subsection{Q7：如何导出或备份课程资料}

\textbf{当前状态}：系统界面暂未提供课程资料的导出功能。

\textbf{替代方案}：

\begin{itemize}
\item 原始文档仍保存在服务器 \texttt{backend/data/uploads/} 目录下，可按目录结构手动获取
\item 知识图谱可通过 Neo4j Browser 导出查询结果
\item 完整备份方法见 7.1 数据备份与恢复
\end{itemize}

\section{学生端学习问题}


\subsection{Q1：加课申请未被通过}

\textbf{问题现象}：提交申请后长时间处于「待审核」。

\textbf{可能原因}：教师尚未处理。

\textbf{解决步骤}：

\begin{enumerate}
\item 联系任课教师，请其在「成员管理」中审核
\item 在课程中心确认申请状态是否为「待审核」
\end{enumerate}
\textbf{仍无法解决时}：请教师核对加课码是否正确，或改用邀请链接。


\subsection{Q2：文档打不开}

\textbf{问题现象}：点击「在线阅读」无反应或报错。

\textbf{可能原因}：文件已被删除；尚未加入课程。

\textbf{解决步骤}：

\begin{enumerate}
\item 确认已成功加入该课程（审核通过）
\item 联系教师确认文档是否仍然存在
\end{enumerate}
\textbf{仍无法解决时}：尝试下载原文件后用本地阅读器打开（若页面提供下载功能）。


\subsection{Q3：知识点无法标记为「已掌握」}

\textbf{问题现象}：点击「已掌握」无反应或提示失败。

\textbf{可能原因}：会话过期；网络异常。

\textbf{解决步骤}：

\begin{enumerate}
\item 刷新页面后重试
\item 若提示登录失效，重新登录
\end{enumerate}
\textbf{仍无法解决时}：确认所选的课程与文档是否正确（标记是针对具体文档下的知识点的）。


\subsection{Q4：学习路径推荐为空}

\textbf{问题现象}：路径推荐页面没有任何推荐结果。

\textbf{可能原因}：未标记任何已掌握知识点；或图谱中没有前置关系。

\textbf{解决步骤}：

\begin{enumerate}
\item 前往「图谱浏览」，标记若干个知识点为「已掌握」
\item 返回「学习路径推荐」刷新查看
\end{enumerate}
\textbf{仍无法解决时}：查看图谱中是否存在前置关系（可请教师在图谱管理中确认）。若图谱没有前置关系，系统会退化为推荐入门知识点。


\subsection{Q5：问答没有引用来源}

\textbf{问题现象}：回答正常，但没有显示引用知识点。

\textbf{可能原因}：未选定课程与文档；或检索结果为空。

\textbf{解决步骤}：

\begin{enumerate}
\item 确认已选择课程与文档（问答是文档级范围的）
\item 确认该文档已完成知识抽取（知识点数 > 0）
\end{enumerate}
\textbf{仍无法解决时}：见 6.5 Q2。


\subsection{Q6：练习提交失败}

\textbf{问题现象}：作答后无法提交。

\textbf{可能原因}：网络异常；会话过期。

\textbf{解决步骤}：

\begin{enumerate}
\item 刷新页面重新登录
\item 重新作答并提交
\end{enumerate}
\textbf{仍无法解决时}：联系教师确认该题目是否仍然有效（已被删除或停用的题目无法提交）。


\section{管理员与权限问题}


\subsection{Q1：如何为忘记密码的用户重置密码}

\textbf{推荐做法}：使用\textbf{管理端的「重置密码」功能}（用户管理页面），而不是直接操作数据库。

\textbf{操作步骤}：

\begin{enumerate}
\item 进入「用户管理」
\item 找到目标用户
\item 执行「重置密码」
\item 将新密码告知用户（用户下次登录需修改）
\end{enumerate}
\textbf{仍无法解决时}：见 7.2 的应急方案说明。


\subsection{Q2：账号被禁用后如何恢复}

\textbf{操作步骤}：进入「用户管理」，找到该账号，执行「启用」。


\subsection{Q3：教师访问其他课程被拒绝}

\textbf{问题现象}：教师点击非本人课程时提示无权限。

\textbf{说明}：\textbf{这是正确的权限行为，不是故障。} 教师对本人课程之外无可见性。公开课程的元数据只对学生开放（学生需要看到详情才能申请加入）。

\textbf{仍无法解决时}：若确需访问，应通过被邀请为「协作教师」的方式加入该课程。


\subsection{Q4：学生误入教师页面}

\textbf{问题现象}：学生点击「数据总览」后跳转到「学习总览」。

\textbf{说明}：\textbf{这是预期行为。} 系统会自动把学生引导到学习总览，避免学生进入教师界面后触发权限报错。


\subsection{Q5：课程如何下架或归档}

\textbf{操作入口}：左侧导航 → 「课程治理」

\textbf{说明}：执行治理操作后，被下架的课程不再对非成员学生暴露元数据。注意「下架」不会改写课程的成员关系。


\subsection{Q6：如何查看系统运行状态}

\textbf{操作入口}：左侧导航 → 「系统监控」

\textbf{说明}：展示请求量、响应耗时与错误数。这些数据来自内存实时计数，\textbf{进程重启后清零}，不提供历史趋势。如需历史数据，应接入外部监控系统。


\section{大模型、向量检索与数据库问题}


\subsection{Q1：知识抽取一直失败}

\textbf{问题现象}：上传文档后抽取阶段持续报错。

\textbf{可能原因}：API Key 无效；余额不足；网络不通；外部服务异常。

\textbf{解决步骤}：

\begin{enumerate}
\item 确认 \texttt{.env} 中 \texttt{LLM\_API\_KEY} 与 \texttt{LLM\_API\_BASE} 正确
\item 登录服务商平台确认余额
\item 测试网络连通性
\item 查看后端终端的完整报错
\end{enumerate}
\textbf{仍无法解决时}：临时切换到其他兼容 OpenAI 接口的模型服务，验证是否为服务商侧问题。


\subsection{Q2：问答回答「知识库中没有相关内容」}

\textbf{问题现象}：提问后系统回答未找到相关内容。

\textbf{可能原因}（按可能性排序）：

\begin{enumerate}
\item \textbf{未选定课程与文档}——问答是文档级范围的
\item \textbf{该文档没有知识点}——抽取失败或结果为空
\item \textbf{Neo4j 中没有图谱数据}——见 Q4
\end{enumerate}
\textbf{解决步骤}：

\begin{enumerate}
\item 确认已选定课程与文档
\item 前往「图谱管理」确认该文档有知识点
\item 在 Neo4j Browser 中执行：
\end{enumerate}
\begin{Verbatim}[breaklines=true,breakanywhere=true,fontsize=\small,frame=single,framesep=4pt,rulecolor=\color{codeframe}]
   MATCH (n:KnowledgePoint) RETURN count(n)
\end{Verbatim}

若结果为 0，说明图谱为空

\textbf{仍无法解决时}：见 Q4。


\subsection{Q3：回答质量突然变差，但不报错}

\textbf{问题现象}：问答功能可用，但回答明显不如之前准确，且系统没有报错。

\textbf{可能原因}：\textbf{向量检索静默降级为关键词检索}。当 \texttt{EMBEDDING\_API\_KEY} 未配置或向量服务调用失败时，系统会自动退回关键词检索——功能仍可用，但质量下降且不报错。

\textbf{解决步骤}：

\begin{enumerate}
\item 检查 \texttt{.env} 中的 \texttt{EMBEDDING\_API\_KEY} 是否配置且有效
\item 确认 SiliconFlow 账户余额
\item 查看后端日志中的警告信息（系统在降级时会记录 warning）
\item 确认网络能访问 \texttt{EMBEDDING\_API\_BASE}
\end{enumerate}
\textbf{仍无法解决时}：临时改用其他提供 embedding 接口的服务。


\subsection{Q4：换机器后图谱为空，问答不工作}

\textbf{问题现象}：在另一台电脑上 \texttt{git clone} 项目并启动后，已有课程还在，但图谱是空的。

\textbf{可能原因}：\textbf{Neo4j 的数据不随代码仓库走。} SQLite 数据库文件（\texttt{app.db}）在项目目录内，会随代码一起获取；但 Neo4j 的数据存储在它自己的数据目录中，\texttt{git clone} 拿不到。

\textbf{解决步骤}（任选其一）：

\begin{enumerate}
\item \textbf{重新上传文档}：在新环境的系统中重新上传课程文档，重新抽取
\item \textbf{从旧环境导出导入}：在旧环境用 Neo4j 的 dump/load 工具迁移数据
\item \textbf{整体迁移}：把旧环境的 Neo4j 数据目录整体复制到新环境
\end{enumerate}
\textbf{仍无法解决时}：确认新环境的 Neo4j 确实在运行且 \texttt{.env} 指向正确的数据库。


\subsection{Q5：修改知识点后第一次提问特别慢}

\textbf{问题现象}：教师修改了知识点描述，之后第一次提问耗时明显偏长。

\textbf{说明}：\textbf{这是预期行为。} 向量索引的失效粒度是\textbf{整个文档}——修改任意一个知识点会导致该文档全部向量被删除，下次问答时全量重算。详见 4.4.4。

\textbf{解决步骤}：无需处理，等待索引重建完成即可。若希望避免，可在修改后先提一个无关问题把索引「热」起来。

\textbf{仍无法解决时}：若耗时异常长，确认该文档的知识点数量是否过大。


\subsection{Q6：数据库文件损坏或数据异常}

\textbf{问题现象}：SQLite 报错，或数据出现异常。

\textbf{解决步骤}：

\begin{enumerate}
\item \textbf{立即停止后端服务}，避免继续写入
\item 备份当前的 \texttt{backend/data/app.db}（即使已损坏，也保留以便分析）
\item 从最近的备份恢复（见 7.1）
\end{enumerate}
\textbf{仍无法解决时}：SQLite 自带完整性检查，可执行 \texttt{PRAGMA integrity\_check;} 查看损坏程度。


\subsection{Q7：测试脚本运行后演示数据被污染}

\textbf{问题现象}：运行了 \texttt{test\_*.py} 或前端 e2e 测试后，演示库中多出了测试账号或课程。

\textbf{可能原因}：\textbf{多数验证脚本会真实写入数据库}——会真实创建账号、课程与节点，而不是使用隔离的测试库。

\textbf{解决步骤}：

\begin{enumerate}
\item 停止使用演示库运行测试
\item 手动清理测试数据，或从备份恢复
\item \textbf{今后在运行测试前先备份 \texttt{app.db}}，或另建一份数据库专用于测试
\end{enumerate}
\begin{todobox}
\warnmark{} \textbf{这是本项目当前的一个已知风险}：验证脚本默认连接线上库，没有独立的测试环境配置。建议在演示前统一执行一次数据清理，并核对 \texttt{t\_course}、\texttt{t\_document}、\texttt{t\_question}、\texttt{t\_kp\_embedding} 等关键表的计数。

\end{todobox}


\section{安全与凭据问题}


\subsection{Q1：忘记引导账号的密码}

\textbf{推荐做法}：优先使用管理端的「重置密码」功能（需另一个可用的管理员账号）。

\textbf{备选做法}：查看后端启动日志——随机密码在账号首次创建时打印过一次。日志文件位于 \texttt{backend/\_uvicorn.log}，搜索「管理员」或「演示教师」。

\textbf{应急做法}（仅限开发环境，见 7.2）：删除该账号后重启后端重建。


\subsection{Q2：怀疑凭据泄露}

\textbf{处理步骤}：

\begin{enumerate}
\item \textbf{更换 \texttt{SECRET\_KEY}}（会使全部已签发的登录凭证失效，所有用户需重新登录）
\item 重置相关账号密码
\item 检查审计日志（管理端 → 审计日志）确认是否有异常操作
\item 按 7.3 处理可能泄露密码的日志文件
\end{enumerate}

\subsection{Q3：日志中是否有敏感信息}

\textbf{说明}：当引导账号使用\textbf{随机密码}创建时，该密码会被打印到启动日志。这是为了让部署者能够获取初始凭据，但也意味着\textbf{日志文件中可能包含有效凭据}。

\textbf{处理建议}：见 7.3 日志及系统监控。


\chapter{数据管理与系统维护}


\section{数据备份与恢复}


\subsection{需要备份的数据}

系统的数据分布在\textbf{三个位置}，\textbf{缺一不可}：

\begin{center}
\small
\begin{tabularx}{\linewidth}{@{}XXXX@{}}
\toprule
\# & 数据 & 位置 & 内容 \\
\midrule
1 & SQLite 数据库 & \texttt{backend/data/app.db} & 用户、课程、文档元数据、题库、学习记录、全部向量 \\
2 & 上传的原始文件 & \texttt{backend/data/uploads/} & 按 \texttt{\{课程ID\}/\{文档ID\}\_\{文件名\}} 组织 \\
3 & Neo4j 数据 & Neo4j 的 \texttt{data/} 目录 & 知识图谱本体（节点与关系） \\
\bottomrule
\end{tabularx}
\end{center}
\begin{todobox}
\warnmark{} \textbf{只备份其中一项是不够的}：SQLite 里没有图数据，Neo4j 里没有用户和向量，uploads 里只有原始文件。\textbf{三者缺一，系统都无法恢复到可用状态。}

\end{todobox}


\subsection{备份方法}

\textbf{SQLite 与上传文件}：

\begin{enumerate}
\item \textbf{先停止后端服务}（避免备份过程中数据正在写入）
\item 直接复制 \texttt{backend/data/} 整个目录
\end{enumerate}
\begin{notebox}
\textbf{为什么建议停服}：SQLite 采用 WAL 模式，运行中直接复制文件可能出现主库文件与 WAL 文件不一致的情况。停服是最稳妥的方式。

\end{notebox}

\textbf{Neo4j}：

\begin{itemize}
\item \textbf{Neo4j Desktop}：使用其内置的 Dump / Backup 功能
\item \textbf{命令行}：停库后复制整个数据目录，或使用 \texttt{neo4j-admin database dump}
\end{itemize}

\subsection{恢复流程（含顺序要求）}

\textbf{必须按以下顺序恢复}，否则可能出现「有文档记录但文件不存在」或「有用户但图谱为空」的不一致状态：

\begin{center}
\small
\begin{tabularx}{\linewidth}{@{}XXX@{}}
\toprule
步骤 & 操作 & 说明 \\
\midrule
1 & 停止后端服务 & 避免恢复期间有写入 \\
2 & 停止 Neo4j &  \\
3 & 恢复 \textbf{SQLite}：将备份的 \texttt{app.db} 覆盖到 \texttt{backend/data/app.db} &  \\
4 & 恢复 \textbf{上传文件}：将备份的 \texttt{uploads/} 覆盖到 \texttt{backend/data/uploads/} &  \\
5 & 恢复 \textbf{Neo4j 数据}：通过 Dump 载入或覆盖数据目录 &  \\
6 & 启动 Neo4j，确认状态正常 &  \\
7 & 启动后端，观察启动日志是否有迁移报错 &  \\
8 & 执行恢复后验证（见下表） &  \\
\bottomrule
\end{tabularx}
\end{center}

\subsection{恢复后的一致性验证}

\begin{center}
\small
\begin{tabularx}{\linewidth}{@{}XXXX@{}}
\toprule
\# & 验证项 & 方法 & 期望 \\
\midrule
1 & 后端启动无异常 & 观察启动日志 & 无迁移报错 \\
2 & 健康接口 & \texttt{curl http://localhost:8000/health} & 返回 \texttt{\{"status":"ok"\}} \\
3 & 图谱数据完整 & Neo4j Browser：\texttt{MATCH (n:KnowledgePoint) RETURN count(n)} & 与备份前一致 \\
4 & 文档记录与文件一致 & 打开若干文档的「在线阅读」 & 均能正常打开 \\
5 & 用户可登录 & 用备份前的账号登录 & 登录成功 \\
6 & 问答可用 & 对已抽取的文档提问 & 返回回答与引用来源 \\
\bottomrule
\end{tabularx}
\end{center}
\begin{todobox}
\warnmark{} \textbf{版本兼容性}：从\textbf{较新版本的 Neo4j 备份}恢复到\textbf{较旧版本}通常不可行（Neo4j 不支持降级恢复）。恢复前请确认目标环境的 Neo4j 版本不低于备份时的版本。

\end{todobox}


\section{账号与密钥管理}


\subsection{密码重置（推荐方式）}

\textbf{优先使用管理端功能}：

\begin{enumerate}
\item 进入「用户管理」
\item 找到目标用户
\item 执行「重置密码」
\item 将新密码告知用户（用户下次登录会被要求修改）
\end{enumerate}
\textbf{优点}：不直接操作数据库，不影响关联数据，且会触发「强制修改密码」流程。


\subsection{密码重置（应急方式，仅限开发环境）}

\textbf{仅在以下情况使用}：没有任何可用的管理员账号，且已确认无法通过其他途径恢复。

\begin{todobox}
\warnmark{} \textbf{直接删除数据库中的用户记录会带来副作用}：

\begin{itemize}
\item 该用户的学习记录、收藏、作答记录等\textbf{引用该用户 ID 的数据会变成孤立数据}；
\item 与审计日志中记录的该用户操作\textbf{失去对应关系}；
\item 若该用户是课程创建者，课程将失去归属。
\end{itemize}
\textbf{操作前必须完整备份 \texttt{app.db}。}

\end{todobox}

\textbf{步骤}（开发环境应急）：

\begin{enumerate}
\item 停止后端服务
\item 备份 \texttt{backend/data/app.db}
\item 确认该账号确实需要重建，并记录其 \texttt{user\_id} 以便后续清理孤立数据
\item 通过 SQLite 客户端删除该用户记录
\item 配置 \texttt{DEFAULT\_ADMIN\_PASSWORD} 或 \texttt{DEFAULT\_TEACHER\_PASSWORD} 环境变量
\item 重启后端，账号将以新密码重建
\end{enumerate}
\begin{notebox}
\textbf{注意}：环境变量\textbf{仅在账号不存在时生效}，不会覆盖已有账号的密码。

\end{notebox}


\subsection{密钥管理}

\begin{center}
\small
\begin{tabularx}{\linewidth}{@{}XXX@{}}
\toprule
密钥 & 存放位置 & 轮换影响 \\
\midrule
\texttt{SECRET\_KEY} & \texttt{backend/.env} & 更换后\textbf{全部登录凭证失效}，所有用户需重新登录 \\
\texttt{LLM\_API\_KEY} & \texttt{backend/.env} & 无影响，直接替换 \\
\texttt{EMBEDDING\_API\_KEY} & \texttt{backend/.env} & 无影响，直接替换 \\
\texttt{NEO4J\_PASSWORD} & \texttt{backend/.env} & 需同时在 Neo4j 侧修改 \\
\bottomrule
\end{tabularx}
\end{center}
\textbf{建议}：

\begin{itemize}
\item \texttt{.env} \textbf{不纳入版本控制}，限制文件权限为仅运行账号可读
\item 定期轮换 API Key
\item 演示环境与生产环境使用不同的密钥
\end{itemize}

\section{日志及系统监控}


\subsection{日志中的敏感信息}

\begin{center}
\small
\begin{tabularx}{\linewidth}{@{}XX@{}}
\toprule
场景 & 是否包含敏感信息 \\
\midrule
引导账号使用\textbf{随机密码}创建 & \warnmark{} \textbf{是}——密码会打印到启动日志，且仅打印一次 \\
引导账号使用\textbf{环境变量密码}创建 & 否 \\
运行期常规日志 & 通常不含凭据 \\
\bottomrule
\end{tabularx}
\end{center}
\textbf{处理建议}：

\begin{center}
\small
\begin{tabularx}{\linewidth}{@{}XX@{}}
\toprule
措施 & 说明 \\
\midrule
访问控制 & 限制日志目录的读取权限，仅运维人员可访问 \\
保留期限 & 保留必要期限后删除，不建议长期留存 \\
安全删除 & 确认不再需要后，彻底删除含初始密码的日志 \\
及时改密 & 首次登录后立即修改密码（系统会强制要求），使日志中的密码失效 \\
\bottomrule
\end{tabularx}
\end{center}
\begin{notebox}
\textbf{最根本的解决办法}：正式部署时显式配置 \texttt{DEFAULT\_ADMIN\_PASSWORD} 与 \texttt{DEFAULT\_TEACHER\_PASSWORD}，从源头避免密码被写入日志。

\end{notebox}


\subsection{系统监控}

\textbf{入口}：管理端 → 「系统监控」

\textbf{指标}：请求量、响应耗时、错误数。

\textbf{重要说明}：这些指标来自\textbf{内存中的实时计数}，不写数据库、不落日志。这样设计的目的是避免监控本身成为新的性能负担或故障点。\textbf{因此进程重启后计数清零}，它反映的是本次运行期间的情况。

\textbf{如果需要长期监控}：应接入外部监控系统（如 Prometheus），从应用日志或新增的指标接口采集。


\section{数据清理与版本升级}


\subsection{删除操作的影响范围}

删除课程或文档时，系统会执行\textbf{级联清理}，同时清理两套数据库中的相关数据。

\textbf{删除文档时清理的内容}：

\begin{Verbatim}[breaklines=true,breakanywhere=true,fontsize=\small,frame=single,framesep=4pt,rulecolor=\color{codeframe}]
(1) Neo4j 中该文档的知识图谱（节点与关系）
(2) SQLite 中该文档的向量
(3) 该文档相关的学习记录
(4) 该文档相关的收藏
(5) 该文档的题目及其作答记录、题目收藏
(6) 服务器上的原始文件
(7) 文档记录本身
\end{Verbatim}

\textbf{两点需要注意}：

\begin{enumerate}
\item \textbf{学生的相关学习记录会被一并删除}——删除文档不只是删掉一份文件，还包括学生在这份文档上的学习痕迹。
\item \textbf{课程通用题会被保留}——不属于任何文档的题目（\texttt{document\_id} 为空）不会随文档删除而消失。
\end{enumerate}

\subsection{版本升级注意事项}

\begin{center}
\small
\begin{tabularx}{\linewidth}{@{}XX@{}}
\toprule
事项 & 说明 \\
\midrule
升级前备份 & \textbf{必须先完整备份}（见 7.1） \\
数据库迁移 & 系统启动时会自动执行幂等的表结构迁移，无需手工操作 \\
依赖版本 & 不要单独升级 \texttt{starlette}（见 6.1 Q2） \\
回滚 & 若升级失败，用备份回滚；注意 Neo4j 不支持降级恢复 \\
\bottomrule
\end{tabularx}
\end{center}

\section{故障应急处理}


\subsection{服务无法启动}

\begin{center}
\small
\begin{tabularx}{\linewidth}{@{}XX@{}}
\toprule
步骤 & 操作 \\
\midrule
1 & 确认 Neo4j 已启动 \\
2 & 确认 \texttt{.env} 配置完整 \\
3 & 查看后端终端报错 \\
4 & 参照 6.1 定位具体问题 \\
5 & 若无法解决，从最近备份恢复 \\
\bottomrule
\end{tabularx}
\end{center}

\subsection{外部服务不可用}

\begin{center}
\small
\begin{tabularx}{\linewidth}{@{}XXX@{}}
\toprule
服务 & 影响 & 应对 \\
\midrule
DeepSeek 不可用 & 无法抽取新文档；问答无法生成回答 & 等待恢复，或临时切换到其他兼容 OpenAI 接口的模型 \\
向量服务不可用 & 问答\textbf{自动降级}为关键词检索，功能可用但质量下降 & 检查 key 与余额；恢复后下次问答会自动重建索引 \\
\bottomrule
\end{tabularx}
\end{center}
\begin{notebox}
\textbf{降级是静默的}，不会报错。排查「回答质量变差」时请优先检查此处（见 6.5 Q3）。

\end{notebox}


\subsection{演示前检查清单}

\begin{center}
\small
\begin{tabularx}{\linewidth}{@{}XXXX@{}}
\toprule
\# & 检查项 & 命令/操作 & 通过 \\
\midrule
1 & Neo4j 运行中 & 访问 \texttt{http://localhost:7474} & $\square$ \\
2 & 后端健康 & \texttt{curl http://localhost:8000/health} & $\square$ \\
3 & 前端可访问 & 打开 \texttt{http://localhost:5173} & $\square$ \\
4 & 图谱非空 & Neo4j：\texttt{MATCH (n:KnowledgePoint) RETURN count(n)} & $\square$ \\
5 & 演示文档已抽取完成 & 课程文档列表显示知识点数 > 0 & $\square$ \\
6 & \textbf{索引已热身} & 提前对演示文档提一个问题 & $\square$ \\
7 & 演示账号可登录 & 教师账号与学生账号各登录一次 & $\square$ \\
8 & 清理测试数据 & 核对课程/文档/题目计数 & $\square$ \\
9 & 网络通畅 & 能访问 DeepSeek 与 SiliconFlow & $\square$ \\
\bottomrule
\end{tabularx}
\end{center}
\begin{notebox}
\textbf{第 6 项尤其重要}：提前提问可以把向量索引构建好，避免现场演示时第一次问答耗时过长（见 4.4.4）。

\end{notebox}


\chapter*{附录}
\addcontentsline{toc}{chapter}{附录}


\section*{附录 A 系统目录结构}

见 2.3。


\section*{附录 B 环境变量速查}

\begin{center}
\small
\begin{tabularx}{\linewidth}{@{}XXXX@{}}
\toprule
变量 & 必填 & 默认值 & 说明 \\
\midrule
\texttt{LLM\_API\_KEY} & $\checkmark$ & — & DeepSeek API Key \\
\texttt{LLM\_API\_BASE} & $\checkmark$ & \texttt{https://api.deepseek.com/v1} & LLM 接口地址 \\
\texttt{LLM\_MODEL} & $\checkmark$ & \texttt{deepseek-chat} & 模型名 \\
\texttt{EMBEDDING\_API\_KEY} & $\checkmark$ & — & 缺失则 RAG 降级为关键词检索 \\
\texttt{EMBEDDING\_API\_BASE} & $\checkmark$ & \texttt{https://api.siliconflow.cn/v1} & 向量接口地址 \\
\texttt{EMBEDDING\_MODEL} & $\checkmark$ & \texttt{BAAI/bge-m3} & 向量模型 \\
\texttt{NEO4J\_URI} & $\checkmark$ & \texttt{bolt://localhost:7687} & Neo4j 连接地址 \\
\texttt{NEO4J\_USER} & $\checkmark$ & \texttt{neo4j} & Neo4j 用户名 \\
\texttt{NEO4J\_PASSWORD} & $\checkmark$ & — & Neo4j 密码 \\
\texttt{SECRET\_KEY} & $\checkmark$ & — & JWT 签名密钥，\textbf{必须修改} \\
\texttt{SQLITE\_DB\_PATH} & $\square$ & \texttt{./data/app.db} & 相对 \texttt{backend/} 目录 \\
\texttt{UPLOAD\_DIR} & $\square$ & \texttt{./data/uploads} & 相对 \texttt{backend/} 目录 \\
\texttt{DEFAULT\_ADMIN\_USERNAME} & $\square$ & \texttt{sysadmin} & 管理员引导账号名 \\
\texttt{DEFAULT\_ADMIN\_PASSWORD} & $\square$ & 随机生成 & \textbf{仅账号不存在时生效} \\
\texttt{DEFAULT\_TEACHER\_PASSWORD} & $\square$ & 随机生成 & \textbf{仅账号不存在时生效} \\
\bottomrule
\end{tabularx}
\end{center}

\section*{附录 C 接口速查}

\begin{notebox}
本附录面向部署与联调人员，普通用户无需阅读。

\end{notebox}

\textbf{认证要求}：除注册、登录与健康检查外，所有接口均需在请求头携带 JWT 令牌。

\textbf{接口前缀现状}：

\begin{center}
\small
\begin{tabularx}{\linewidth}{@{}XX@{}}
\toprule
模块 & 前缀 \\
\midrule
用户认证 & \texttt{/api/auth} \warnmark{} \\
知识图谱管理 & \texttt{/api/kg} \warnmark{} \\
课程管理 & \texttt{/api/v1/courses} \\
课程成员 & \texttt{/api/v1/courses} \\
课程邀请 & \texttt{/api/v1/invites} \\
文档管理 & \texttt{/api/v1/documents} \\
智能问答 & \texttt{/api/v1/qa} \\
学习路径 & \texttt{/api/v1/learning-path} \\
学习记录 & \texttt{/api/v1/learning} \\
知识图谱（关系补全） & \texttt{/api/v1/graph} \\
数据总览 & \texttt{/api/v1/dashboard} \\
收藏夹 & \texttt{/api/v1/favorites} \\
教师监测 & \texttt{/api/v1/teacher} \\
题库管理 & \texttt{/api/v1/questions} \\
做题练习 & \texttt{/api/v1/practice} \\
主观题批改 & \texttt{/api/v1/grading} \\
个人中心 & \texttt{/api/v1/profile} \\
管理员端 & \texttt{/api/v1/admin} \\
\bottomrule
\end{tabularx}
\end{center}
\begin{todobox}
\warnmark{} 标注的 \texttt{/api/auth} 与 \texttt{/api/kg} 两个前缀\textbf{尚未迁移到 \texttt{/api/v1/} 规范下}。前端已按实际路径适配，功能不受影响。建议在后续版本中统一，或在《项目详细方案》中记录为待办。

\textbf{健康检查}：\texttt{GET /health} → \texttt{\{"status":"ok"\}}

\end{todobox}


\section*{附录 D 测试脚本说明}

\begin{todobox}
\warnmark{} \textbf{重要提醒：多数验证脚本会真实写入数据库}（真实创建账号、课程与节点），而不是使用隔离的测试库。\textbf{运行前请先备份 \texttt{app.db}}，或在专用环境运行。

\end{todobox}


\subsection*{端到端验证脚本（\texttt{backend/test\_*.py}）}

自带断言，直接运行，输出形如「通过 N/M」。

\begin{center}
\small
\begin{tabularx}{\linewidth}{@{}XX@{}}
\toprule
脚本 & 用途 \\
\midrule
\texttt{test\_pipeline.py} & 全链路冒烟测试 \\
\texttt{test\_upload\_flow.py} & 上传 → 解析 → 抽取 → 入图 \\
\texttt{test\_path\_recommender.py} & 学习路径推荐 \\
\texttt{test\_vector\_index.py} & 向量索引构建与失效 \\
\texttt{test\_graph\_integrity.py} & 图谱完整性（悬空边等） \\
\texttt{test\_permission\_matrix.py} & 权限矩阵验证 \\
\texttt{test\_perf\_timing.py} & \textbf{性能计时（对应 60 秒/15 秒指标）} \\
\texttt{test\_qa\_testset.py} & 问答测试集 \\
\bottomrule
\end{tabularx}
\end{center}

\subsection*{评测脚本（\texttt{backend/eval\_*.py}）}

只读，结果写入 \texttt{backend/eval\_data/*.json}。

\begin{center}
\small
\begin{tabularx}{\linewidth}{@{}XX@{}}
\toprule
脚本 & 用途 \\
\midrule
\texttt{eval\_accuracy.py} & 知识抽取准确率（P/R/F1，全仓库唯一评测实现） \\
\texttt{eval\_ann\_compare.py} & 近似最近邻算法对比 \\
\texttt{eval\_bloom\_filter.py} & 布隆过滤器评测 \\
\bottomrule
\end{tabularx}
\end{center}

\subsection*{前端 e2e（\texttt{frontend-vue/tests/*.spec.ts}）}

需前后端均已启动。

\begin{Verbatim}[breaklines=true,breakanywhere=true,fontsize=\small,frame=single,framesep=4pt,rulecolor=\color{codeframe}]
cd frontend-vue
npx playwright test
\end{Verbatim}

\begin{todobox}
\warnmark{} e2e 测试同样会向真实库写入测试账号与课程，跑完请清理 \texttt{e2e\_} / \texttt{E2E} 前缀的数据。

\end{todobox}


\section*{附录 E 赛题要求—功能—章节—证据映射表}

\begin{center}
\footnotesize
\begin{tabularx}{\linewidth}{@{}XXXXX@{}}
\toprule
赛题要求 & 系统功能 & 文档章节 & 建议提供的证据 & 状态 \\
\midrule
支持 ≥2 种文档格式 & PDF/DOCX/TXT/MD 四种格式解析 & 3.2.3 & 实际上传与解析成功截图 & \shotmark{} 待补 \\
知识点 ≥20 个（一章内容） & 自动抽取知识点实体 & 3.2.3 5.3 & 真实课程抽取结果截图 & \shotmark{} 待补 \\
关系类型 ≥3 种 & CONTAINS/PRECEDES/APPLIES\_TO/RELATED\_TO & 4.2 5.3 & 关系类型统计与图谱截图 & \shotmark{} 待补 \\
抽取准确率 ≥70\% & LLM 抽取 + 多轮提示词迭代 & 5.3 & 独立测试报告（F1=0.8312） & $\checkmark$ 已有报告 \\
图谱交互与详情展示 & 缩放/拖拽/点击详情/四种布局 & 3.2.4 & 交互操作截图 & \shotmark{} 待补 \\
列表/卡片展示知识点详情 & 知识点详情抽屉 & 3.2.4 & 详情展示截图 & \shotmark{} 待补 \\
智能问答（RAG） & 向量检索 + 大模型生成 + 引用来源 & 4.4 & 测试问题、答案、引用与人工评价 & \todomark{} 待补 \\
学习路径推荐 & 前置关系图遍历 + 优先级排序 & 4.3 & 已掌握状态变化前后的推荐结果对比 & \shotmark{} 待补 \\
路径可视化展示 & 路径链路图 & 3.3.3 & 路径界面截图 & \shotmark{} 待补 \\
教师端功能完整 & 上传/进度/预览/手动修正 & 3.2 & 编辑前后对比截图 & \shotmark{} 待补 \\
学生端功能完整 & 浏览/详情/问答/标记/推荐 & 3.3 & 各功能截图 & \shotmark{} 待补 \\
多课程管理 ≥2 门 & 课程与文档级隔离 & 5.5 & 两门课程的隔离验证记录 & \todomark{} 待补 \\
解析+抽取 ≤60 秒 & 分块并发抽取 & 5.4 & 实际性能测试记录 & \todomark{} 待补 \\
问答 ≤15 秒 & 向量检索 + 生成 & 5.4 & 实际性能测试记录 & \todomark{} 待补 \\
普通电脑可运行 & 无需 GPU & 2.2 & 部署环境说明 & $\checkmark$ \\
\bottomrule
\end{tabularx}
\end{center}
\textbf{配套材料体系}：

\begin{center}
\small
\begin{tabularx}{\linewidth}{@{}XXX@{}}
\toprule
材料 & 负责内容 & 与本文档的关系 \\
\midrule
本文档 & 如何使用、如何验证 & — \\
《知识抽取准确率测试报告》 & 评估方法与实验结果 & 本文档 5.3 引用其结论 \\
《问答测试集》 & 问答的测试题与结果 & 本文档 5.2.1 引用其结论 \\
《提示词工程完整记录》 & 提示词迭代过程 & 本文档 4.1.3 引用其规则 \\
《项目详细方案》 & 架构设计、关键算法、技术选型 & 本文档第四章的原理部分引用其细节 \\
\bottomrule
\end{tabularx}
\end{center}

\section*{附录 F 术语表}

\begin{center}
\small
\begin{tabularx}{\linewidth}{@{}XXX@{}}
\toprule
术语 & 英文/原文 & 通俗解释 \\
\midrule
知识图谱 & Knowledge Graph & 用「节点 + 连线」表示知识及其关联的结构 \\
知识点 & Knowledge Point & 图谱中的一个节点，如「二叉树」「快速排序」 \\
实体 & Entity & 同上，指抽取出的知识点对象 \\
关系 & Relation & 知识点之间的连线，本系统有四种 \\
\texttt{PRECEDES} & 前置关系 & 「先学 A 才能学 B」的依赖关系，是路径推荐的基础 \\
\texttt{CONTAINS} & 包含关系 & 「A 包含 B」，如「线性表包含栈」 \\
\texttt{APPLIES\_TO} & 应用关系 & 「A 用于解决 B」，如「哈希算法应用于查找问题」 \\
\texttt{RELATED\_TO} & 相关关系 & 密切相关但不是以上三类 \\
前置知识 & Prerequisite & 学习某知识点前需要先掌握的知识点 \\
根节点 / 入门知识点 & Root Node & 没有任何前置关系的知识点，是学习的起点 \\
RAG & Retrieval-Augmented Generation & 检索增强生成：先检索资料，再让大模型依据资料作答 \\
向量 / Embedding & Embedding & 把文字转换成一串数字，使语义相近的文字在数值上也接近 \\
余弦相似度 & Cosine Similarity & 衡量两个向量方向接近程度的指标，用于判断语义相关度 \\
索引（向量索引） & Vector Index & 为知识点预先算好的向量集合，问答时用于快速查找 \\
懒构建 & Lazy Build & 不提前做，等到真正需要时才做（本系统用于向量索引） \\
降级 & Graceful Degradation & 某个功能不可用时自动切换到较弱但可用的替代方案 \\
抽取 & Extraction & 从文档中识别出知识点及其关系的过程 \\
分块 & Chunking & 把长文档切成较小的段落，以便交给大模型处理 \\
悬空边 & Dangling Edge & 连线的一端指向不存在的节点，属于数据错误 \\
级联删除 & Cascade Delete & 删除主数据时一并清理其关联的从属数据 \\
P95 & 95th Percentile & 把耗时可从小到大排序，95\% 的请求都快于这个值 \\
幂等 & Idempotent & 同一操作执行多次与执行一次的结果相同 \\
ANN & Approximate Nearest Neighbor & 近似最近邻：用精度换速度的向量检索方式，本系统规模下未使用 \\
\bottomrule
\end{tabularx}
\end{center}

\section*{附录 G 文档版本与修订记录}

\begin{center}
\small
\begin{tabularx}{\linewidth}{@{}XX@{}}
\toprule
项 & 内容 \\
\midrule
文档版本 & v2.0 \\
对应系统版本 & \verifymark{} 待填（见 \texttt{backend/app/core/config.py} 的 \texttt{VERSION}） \\
发布日期 & \verifymark{} 待填 \\
编制人 & \verifymark{} 待填 \\
修订记录 & 见本文档开头 \\
\bottomrule
\end{tabularx}
\end{center}

\section*{提交前检查清单}

\begin{center}
\small
\begin{tabularx}{\linewidth}{@{}XXX@{}}
\toprule
\# & 检查项 & 完成 \\
\midrule
1 & 所有 \shotmark{} 截图占位符已替换为\textbf{真实系统截图} & $\square$ \\
2 & 所有 \todomark{} 数值占位符已替换为\textbf{实测数据} & $\square$ \\
3 & 所有 \verifymark{} 待确认项已填写（版本号、日期、仓库地址等） & $\square$ \\
4 & 全文搜索「待填」「待补」「\todomark{}」「\shotmark{}」「\verifymark{}」，确保无遗留 & $\square$ \\
5 & 文档中引用数值与《知识抽取准确率测试报告》《问答测试集》\textbf{保持一致} & $\square$ \\
6 & 已在专用环境或先备份后运行测试脚本，演示库未被污染 & $\square$ \\
7 & 已补充 PRECEDES 关系类型的准确率测试数据（见 2.8） & $\square$ \\
8 & 已补充人工抽样验证结果（赛题口径要求） & $\square$ \\
9 & 附录 E 映射表中所有 \todomark{} 项已补齐证据 & $\square$ \\
\bottomrule
\end{tabularx}
\end{center}
\end{document}
